% !TEX program = lualatex
% ---------------------------------------------------------------------------------------------
%  Readback note 5: The parameters and the CM field: places, Frobenius elements, the different
%
%  Build:  python3 tools/rb.py note <slug>
%
%  One unit of a Lean artifact, rewritten as ordinary mathematics. The prose is hand-written. Lean names,
%  line ranges and counts come from data.tex; quoted Lean statements are copied from the artifact by the build.
%
%  Rules of the text: standard terminology; every symbol is explained where it first appears and listed again
%  in the notation at the end; no letter has two meanings; a proof is complete or says what it takes for granted;
%  Lean names appear only in the LEAN lines and in the concordance.
%
%  The signature rule (the build enforces it): a function, relation or predicate that the artifact defines is
%  written only through its notation macro, which the build makes from the spec (`notation`) with one slot for
%  each explicit argument of the Lean definition. No argument is fixed and dropped, also not "for brevity" and
%  not for one section. Every argument and every hypothesis of a stated declaration stands inside its
%  statement, or in a Setting that the statement cites by number; only implicit and instance arguments (the
%  ambient objects) may stand in the conventions.
% ---------------------------------------------------------------------------------------------
\documentclass[11pt,a4paper]{article}
\ifdefined\pdfvariable\pdfvariable suppressoptionalinfo \numexpr 1+2+4+8\relax\fi
\usepackage[a4paper,hmargin=35mm,top=22mm,bottom=27mm,footskip=12mm,nohead]{geometry}

\usepackage{fontspec}
\directlua{
  luaotfload.add_fallback("gfallback", {"STIXTwoMath-Regular.otf:mode=node;", "DejaVuSans.ttf:mode=node;"})
  luaotfload.add_fallback("gmonofallback", {"DejaVuSansMono.ttf:mode=node;", "STIXTwoMath-Regular.otf:mode=node;"})
}
\setmainfont{NewCM10}[Extension=.otf,UprightFont=*-Book,ItalicFont=*-BookItalic,
  BoldFont=*-Bold,BoldItalicFont=*-BoldItalic,SmallCapsFont=*-Book,SmallCapsFeatures={Letters=SmallCaps},
  RawFeature={fallback=gfallback}]
\setsansfont{NewCMSans10}[Extension=.otf,UprightFont=*-Book,ItalicFont=*-BookOblique,
  BoldFont=*-Bold,BoldItalicFont=*-BoldOblique,RawFeature={fallback=gfallback}]
\setmonofont{JetBrainsMono}[Extension=.otf,UprightFont=*-Regular,ItalicFont=*-Italic,
  BoldFont=*-SemiBold,BoldItalicFont=*-SemiBoldItalic,Scale=MatchLowercase,
  RawFeature={-calt;-liga;fallback=gmonofallback}]
\usepackage{amsmath}
\usepackage{amsthm}
\usepackage[math-style=TeX]{unicode-math}
\setmathfont{NewCMMath-Book.otf}
\emergencystretch=2.5em\relax

\usepackage{xcolor}
\usepackage{tikz}
\usetikzlibrary{calc,arrows.meta,positioning,backgrounds}
\usepackage[protrusion=true,expansion=true]{microtype}
\usepackage{enumitem}
\usepackage{needspace}
\usepackage{booktabs}
\usepackage{array}
\usepackage{longtable}
\usepackage{fancyhdr}
\usepackage{lastpage}
\usepackage{multicol}
\usepackage[hidelinks,pdfusetitle]{hyperref}

\definecolor{ink}{HTML}{0F172A}
\definecolor{inksoft}{HTML}{475569}
\definecolor{inkfaint}{HTML}{94A3B8}
\definecolor{rulegrey}{HTML}{CBD5E1}
\definecolor{arith}{HTML}{C2410C}
\definecolor{est}{HTML}{7C3AED}
\definecolor{ana}{HTML}{0891B2}
\definecolor{ok}{HTML}{047857}
\definecolor{leancol}{HTML}{3730A3}

\title{The parameters and the CM field: places, Frobenius elements, the different}
\author{grh-compact}
\date{}

% ------------------------------------------------------------------------------------- macros
\newcommand{\leanfont}{\ttfamily\fontsize{8.6}{10.5}\selectfont\color{leancol}}
\newcommand{\lean}[1]{\csname LN@#1\endcsname}
\newcommand{\fillin}[1]{{\color{arith}\itshape[#1]}}  % a place of the template that is not written yet; the build reports every one that is left
\newcommand{\leanquote}[1]{\csname LQ@#1\endcsname}
% generated by rb.py note; do not edit
\newcommand{\ConcDecls}{96}
\newcommand{\Pin}{72fead10}
\newcommand{\PinText}{that of commit {\leanfont 72fead10}}
\newcommand{\ArtifactName}{grh-compact}
\expandafter\newcommand\csname LN@par\endcsname{\href{https://github.com/lsernau/grh-lean/blob/72fead102f663d5c7a407964a13dc0cf9d977207/grh-compact/src/New/Data/Different.lean\#L23-L28}{{\leanfont Parameters}}}
\expandafter\newcommand\csname LQ@par\endcsname{\leanquoteline{0}{@[local grind] structure Parameters where}\leanquoteline{2}{(N₀ d : ℕ)}\leanquoteline{2}{dvd\_N₀ : 18480 ∣ N₀}\leanquoteline{2}{d\_prime : d.Prime}\leanquoteline{2}{/-- `d ≡ N₀/2 - 1 mod N₀`. -/}\leanquoteline{2}{d\_congr : N₀ ∣ 2 * (d + 1) ∧ ¬N₀ ∣ d + 1}}
\expandafter\newcommand\csname LN@twd\endcsname{\href{https://github.com/lsernau/grh-lean/blob/72fead102f663d5c7a407964a13dc0cf9d977207/grh-compact/src/New/Data/Different.lean\#L41-L41}{{\leanfont towerDeg}}}
\expandafter\newcommand\csname LN@Nn\endcsname{\href{https://github.com/lsernau/grh-lean/blob/72fead102f663d5c7a407964a13dc0cf9d977207/grh-compact/src/New/Data/Different.lean\#L44-L44}{{\leanfont Parameters.\allowbreak{}N₁}}}
\expandafter\newcommand\csname LN@KP\endcsname{\href{https://github.com/lsernau/grh-lean/blob/72fead102f663d5c7a407964a13dc0cf9d977207/grh-compact/src/New/Data/Different.lean\#L64-L64}{{\leanfont Parameters.\allowbreak{}K}}}
\expandafter\newcommand\csname LN@Mom\endcsname{\href{https://github.com/lsernau/grh-lean/blob/72fead102f663d5c7a407964a13dc0cf9d977207/grh-compact/src/New/Data/Different.lean\#L68-L68}{{\leanfont Parameters.\allowbreak{}Mω}}}
\expandafter\newcommand\csname LN@heeg\endcsname{\href{https://github.com/lsernau/grh-lean/blob/72fead102f663d5c7a407964a13dc0cf9d977207/grh-compact/src/New/Data/Different.lean\#L36-L37}{{\leanfont Parameters.\allowbreak{}instHeegD}}}
\expandafter\newcommand\csname LN@ltd\endcsname{\href{https://github.com/lsernau/grh-lean/blob/72fead102f663d5c7a407964a13dc0cf9d977207/grh-compact/src/New/Data/Different.lean\#L52-L52}{{\leanfont Parameters.\allowbreak{}lt\_d}}}
\expandafter\newcommand\csname LN@dvdnorm\endcsname{\href{https://github.com/lsernau/grh-lean/blob/72fead102f663d5c7a407964a13dc0cf9d977207/grh-compact/src/New/Data/Different.lean\#L47-L50}{{\leanfont Parameters.\allowbreak{}dvd\_norm\_c₀}}}
\expandafter\newcommand\csname LN@chione\endcsname{\href{https://github.com/lsernau/grh-lean/blob/72fead102f663d5c7a407964a13dc0cf9d977207/grh-compact/src/New/Data/Different.lean\#L55-L58}{{\leanfont Parameters.\allowbreak{}chiK\_eq\_one\_of\_dvd}}}
\expandafter\newcommand\csname LN@thirteen\endcsname{\href{https://github.com/lsernau/grh-lean/blob/72fead102f663d5c7a407964a13dc0cf9d977207/grh-compact/src/New/Data/Different.lean\#L60-L61}{{\leanfont Parameters.\allowbreak{}thirteen\_le\_of\_chiK\_ne\_one}}}
\expandafter\newcommand\csname LN@NdM\endcsname{\href{https://github.com/lsernau/grh-lean/blob/72fead102f663d5c7a407964a13dc0cf9d977207/grh-compact/src/New/Data/Different.lean\#L71-L72}{{\leanfont Parameters.\allowbreak{}N₁\_mul\_d\_dvd\_M\_omega}}}
\expandafter\newcommand\csname LN@coprime\endcsname{\href{https://github.com/lsernau/grh-lean/blob/72fead102f663d5c7a407964a13dc0cf9d977207/grh-compact/src/New/Data/Different.lean\#L74-L75}{{\leanfont Parameters.\allowbreak{}coprime\_N₁\_d}}}
\expandafter\newcommand\csname LN@twopow\endcsname{\href{https://github.com/lsernau/grh-lean/blob/72fead102f663d5c7a407964a13dc0cf9d977207/grh-compact/src/New/Data/Different.lean\#L193-L194}{{\leanfont Parameters.\allowbreak{}two\_pow\_dvd\_N₁}}}
\expandafter\newcommand\csname LN@powd\endcsname{\href{https://github.com/lsernau/grh-lean/blob/72fead102f663d5c7a407964a13dc0cf9d977207/grh-compact/src/New/Data/Different.lean\#L212-L215}{{\leanfont Parameters.\allowbreak{}pow\_d\_ne\_neg\_one}}}
\expandafter\newcommand\csname LN@EP\endcsname{\href{https://github.com/lsernau/grh-lean/blob/72fead102f663d5c7a407964a13dc0cf9d977207/grh-compact/src/New/Data/Different.lean\#L79-L79}{{\leanfont Parameters.\allowbreak{}E}}}
\expandafter\newcommand\csname LN@PlP\endcsname{\href{https://github.com/lsernau/grh-lean/blob/72fead102f663d5c7a407964a13dc0cf9d977207/grh-compact/src/New/Data/Different.lean\#L83-L83}{{\leanfont Parameters.\allowbreak{}Φ}}}
\expandafter\newcommand\csname LN@rtu\endcsname{\href{https://github.com/lsernau/grh-lean/blob/72fead102f663d5c7a407964a13dc0cf9d977207/grh-compact/src/New/Data/Different.lean\#L142-L142}{{\leanfont Parameters.\allowbreak{}ζ}}}
\expandafter\newcommand\csname LN@OmP\endcsname{\href{https://github.com/lsernau/grh-lean/blob/72fead102f663d5c7a407964a13dc0cf9d977207/grh-compact/src/New/Data/Different.lean\#L87-L87}{{\leanfont Parameters.\allowbreak{}bigOmega}}}
\expandafter\newcommand\csname LN@primrt\endcsname{\href{https://github.com/lsernau/grh-lean/blob/72fead102f663d5c7a407964a13dc0cf9d977207/grh-compact/src/New/Data/Different.lean\#L180-L181}{{\leanfont Parameters.\allowbreak{}isPrimitiveRoot\_zeta}}}
\expandafter\newcommand\csname LN@irr\endcsname{\href{https://github.com/lsernau/grh-lean/blob/72fead102f663d5c7a407964a13dc0cf9d977207/grh-compact/src/New/Data/Different.lean\#L92-L100}{{\leanfont Parameters.\allowbreak{}irreducible\_cyclotomic}}}
\expandafter\newcommand\csname LN@tocyc\endcsname{\href{https://github.com/lsernau/grh-lean/blob/72fead102f663d5c7a407964a13dc0cf9d977207/grh-compact/src/New/Data/Different.lean\#L104-L108}{{\leanfont Parameters.\allowbreak{}toCycloField}}}
\expandafter\newcommand\csname LN@embeq\endcsname{\href{https://github.com/lsernau/grh-lean/blob/72fead102f663d5c7a407964a13dc0cf9d977207/grh-compact/src/New/Data/Different.lean\#L116-L117}{{\leanfont Parameters.\allowbreak{}algHom\_eq\_of\_mk\_eq}}}
\expandafter\newcommand\csname LN@galP\endcsname{\href{https://github.com/lsernau/grh-lean/blob/72fead102f663d5c7a407964a13dc0cf9d977207/grh-compact/src/New/Data/Different.lean\#L127-L127}{{\leanfont Parameters.\allowbreak{}galOfPlace}}}
\expandafter\newcommand\csname LN@cycP\endcsname{\href{https://github.com/lsernau/grh-lean/blob/72fead102f663d5c7a407964a13dc0cf9d977207/grh-compact/src/New/Data/Different.lean\#L131-L131}{{\leanfont Parameters.\allowbreak{}placeMulEquiv}}}
\expandafter\newcommand\csname LN@pin\endcsname{\href{https://github.com/lsernau/grh-lean/blob/72fead102f663d5c7a407964a13dc0cf9d977207/grh-compact/src/New/Data/Different.lean\#L135-L135}{{\leanfont Parameters.\allowbreak{}placeOfInt}}}
\expandafter\newcommand\csname LN@pcl\endcsname{\href{https://github.com/lsernau/grh-lean/blob/72fead102f663d5c7a407964a13dc0cf9d977207/grh-compact/src/New/Data/Different.lean\#L139-L139}{{\leanfont Parameters.\allowbreak{}placeClass}}}
\expandafter\newcommand\csname LN@phialg\endcsname{\href{https://github.com/lsernau/grh-lean/blob/72fead102f663d5c7a407964a13dc0cf9d977207/grh-compact/src/New/Data/Different.lean\#L148-L148}{{\leanfont Parameters.\allowbreak{}phi\_algebraMap}}}
\expandafter\newcommand\csname LN@phimul\endcsname{\href{https://github.com/lsernau/grh-lean/blob/72fead102f663d5c7a407964a13dc0cf9d977207/grh-compact/src/New/Data/Different.lean\#L183-L183}{{\leanfont Parameters.\allowbreak{}phi\_mul}}}
\expandafter\newcommand\csname LN@starrt\endcsname{\href{https://github.com/lsernau/grh-lean/blob/72fead102f663d5c7a407964a13dc0cf9d977207/grh-compact/src/New/Data/Different.lean\#L185-L187}{{\leanfont Parameters.\allowbreak{}star\_ζ}}}
\expandafter\newcommand\csname LN@galrt\endcsname{\href{https://github.com/lsernau/grh-lean/blob/72fead102f663d5c7a407964a13dc0cf9d977207/grh-compact/src/New/Data/Different.lean\#L189-L190}{{\leanfont Parameters.\allowbreak{}galOfPlace\_zeta}}}
\expandafter\newcommand\csname LN@pinpow\endcsname{\href{https://github.com/lsernau/grh-lean/blob/72fead102f663d5c7a407964a13dc0cf9d977207/grh-compact/src/New/Data/Different.lean\#L196-L197}{{\leanfont Parameters.\allowbreak{}placeOfInt\_pow}}}
\expandafter\newcommand\csname LN@clpin\endcsname{\href{https://github.com/lsernau/grh-lean/blob/72fead102f663d5c7a407964a13dc0cf9d977207/grh-compact/src/New/Data/Different.lean\#L199-L202}{{\leanfont Parameters.\allowbreak{}placeClass\_placeOfInt}}}
\expandafter\newcommand\csname LN@galpin\endcsname{\href{https://github.com/lsernau/grh-lean/blob/72fead102f663d5c7a407964a13dc0cf9d977207/grh-compact/src/New/Data/Different.lean\#L206-L209}{{\leanfont Parameters.\allowbreak{}galOfPlace\_placeOfInt\_zeta}}}
\expandafter\newcommand\csname LN@clneg\endcsname{\href{https://github.com/lsernau/grh-lean/blob/72fead102f663d5c7a407964a13dc0cf9d977207/grh-compact/src/New/Data/Different.lean\#L217-L220}{{\leanfont Parameters.\allowbreak{}placeClass\_placeOfInt\_eq\_neg\_one}}}
\expandafter\newcommand\csname LN@expandpow\endcsname{\href{https://github.com/lsernau/grh-lean/blob/72fead102f663d5c7a407964a13dc0cf9d977207/grh-compact/src/New/Data/Different.lean\#L156-L158}{{\leanfont Parameters.\allowbreak{}expand\_pow\_cyclotomic}}}
\expandafter\newcommand\csname LN@expandtw\endcsname{\href{https://github.com/lsernau/grh-lean/blob/72fead102f663d5c7a407964a13dc0cf9d977207/grh-compact/src/New/Data/Different.lean\#L162-L164}{{\leanfont Parameters.\allowbreak{}expand\_towerDeg\_cyclotomic}}}
\expandafter\newcommand\csname LN@cardPl\endcsname{\href{https://github.com/lsernau/grh-lean/blob/72fead102f663d5c7a407964a13dc0cf9d977207/grh-compact/src/New/Data/Different.lean\#L168-L170}{{\leanfont Parameters.\allowbreak{}card\_infinitePlace\_E}}}
\expandafter\newcommand\csname LN@cardcl\endcsname{\href{https://github.com/lsernau/grh-lean/blob/72fead102f663d5c7a407964a13dc0cf9d977207/grh-compact/src/New/Data/Different.lean\#L174-L176}{{\leanfont Parameters.\allowbreak{}card\_placeClass\_eq}}}
\expandafter\newcommand\csname LN@prm\endcsname{\href{https://github.com/lsernau/grh-lean/blob/72fead102f663d5c7a407964a13dc0cf9d977207/grh-compact/src/New/Data/Different.lean\#L224-L224}{{\leanfont Parameters.\allowbreak{}placePerm}}}
\expandafter\newcommand\csname LN@prmspec\endcsname{\href{https://github.com/lsernau/grh-lean/blob/72fead102f663d5c7a407964a13dc0cf9d977207/grh-compact/src/New/Data/Different.lean\#L229-L231}{{\leanfont Parameters.\allowbreak{}placePerm\_spec}}}
\expandafter\newcommand\csname LN@roots\endcsname{\href{https://github.com/lsernau/grh-lean/blob/72fead102f663d5c7a407964a13dc0cf9d977207/grh-compact/src/New/Data/Different.lean\#L239-L247}{{\leanfont Parameters.\allowbreak{}exists\_zeta\_pow\_eq}}}
\expandafter\newcommand\csname LN@prodphi\endcsname{\href{https://github.com/lsernau/grh-lean/blob/72fead102f663d5c7a407964a13dc0cf9d977207/grh-compact/src/New/Data/Different.lean\#L249-L250}{{\leanfont Parameters.\allowbreak{}prod\_phi\_eq\_norm}}}
\expandafter\newcommand\csname LN@rti\endcsname{\href{https://github.com/lsernau/grh-lean/blob/72fead102f663d5c7a407964a13dc0cf9d977207/grh-compact/src/New/Data/Different.lean\#L260-L260}{{\leanfont Parameters.\allowbreak{}zetaInt}}}
\expandafter\newcommand\csname LN@czE\endcsname{\href{https://github.com/lsernau/grh-lean/blob/72fead102f663d5c7a407964a13dc0cf9d977207/grh-compact/src/New/Data/Different.lean\#L268-L268}{{\leanfont Parameters.\allowbreak{}c₀InE}}}
\expandafter\newcommand\csname LN@sqd\endcsname{\href{https://github.com/lsernau/grh-lean/blob/72fead102f663d5c7a407964a13dc0cf9d977207/grh-compact/src/New/Data/Different.lean\#L271-L271}{{\leanfont Parameters.\allowbreak{}sqrtNegDInE}}}
\expandafter\newcommand\csname LN@dgen\endcsname{\href{https://github.com/lsernau/grh-lean/blob/72fead102f663d5c7a407964a13dc0cf9d977207/grh-compact/src/New/Data/Different.lean\#L317-L317}{{\leanfont Parameters.\allowbreak{}relDifferentGen}}}
\expandafter\newcommand\csname LN@coert\endcsname{\href{https://github.com/lsernau/grh-lean/blob/72fead102f663d5c7a407964a13dc0cf9d977207/grh-compact/src/New/Data/Different.lean\#L262-L262}{{\leanfont Parameters.\allowbreak{}coe\_zetaInt}}}
\expandafter\newcommand\csname LN@primrti\endcsname{\href{https://github.com/lsernau/grh-lean/blob/72fead102f663d5c7a407964a13dc0cf9d977207/grh-compact/src/New/Data/Different.lean\#L264-L265}{{\leanfont Parameters.\allowbreak{}isPrimitiveRoot\_zetaInt}}}
\expandafter\newcommand\csname LN@coesq\endcsname{\href{https://github.com/lsernau/grh-lean/blob/72fead102f663d5c7a407964a13dc0cf9d977207/grh-compact/src/New/Data/Different.lean\#L276-L277}{{\leanfont Parameters.\allowbreak{}coe\_sqrtNegDInE}}}
\expandafter\newcommand\csname LN@starsq\endcsname{\href{https://github.com/lsernau/grh-lean/blob/72fead102f663d5c7a407964a13dc0cf9d977207/grh-compact/src/New/Data/Different.lean\#L279-L280}{{\leanfont Parameters.\allowbreak{}star\_sqrtNegDInE}}}
\expandafter\newcommand\csname LN@starrtpow\endcsname{\href{https://github.com/lsernau/grh-lean/blob/72fead102f663d5c7a407964a13dc0cf9d977207/grh-compact/src/New/Data/Different.lean\#L282-L283}{{\leanfont Parameters.\allowbreak{}star\_zetaInt\_pow\_mul}}}
\expandafter\newcommand\csname LN@sqdvd\endcsname{\href{https://github.com/lsernau/grh-lean/blob/72fead102f663d5c7a407964a13dc0cf9d977207/grh-compact/src/New/Data/Different.lean\#L285-L286}{{\leanfont Parameters.\allowbreak{}sqrtNegDInE\_dvd}}}
\expandafter\newcommand\csname LN@czmul\endcsname{\href{https://github.com/lsernau/grh-lean/blob/72fead102f663d5c7a407964a13dc0cf9d977207/grh-compact/src/New/Data/Different.lean\#L288-L289}{{\leanfont Parameters.\allowbreak{}c₀InE\_mul\_one\_sub}}}
\expandafter\newcommand\csname LN@starcz\endcsname{\href{https://github.com/lsernau/grh-lean/blob/72fead102f663d5c7a407964a13dc0cf9d977207/grh-compact/src/New/Data/Different.lean\#L291-L292}{{\leanfont Parameters.\allowbreak{}star\_c₀InE}}}
\expandafter\newcommand\csname LN@relspec\endcsname{\href{https://github.com/lsernau/grh-lean/blob/72fead102f663d5c7a407964a13dc0cf9d977207/grh-compact/src/New/Data/Different.lean\#L299-L313}{{\leanfont Parameters.\allowbreak{}relDifferent\_spec}}}
\expandafter\newcommand\csname LN@dgendvd\endcsname{\href{https://github.com/lsernau/grh-lean/blob/72fead102f663d5c7a407964a13dc0cf9d977207/grh-compact/src/New/Data/Different.lean\#L319-L319}{{\leanfont Parameters.\allowbreak{}relDifferentGen\_dvd\_N₁}}}
\expandafter\newcommand\csname LN@dgenne\endcsname{\href{https://github.com/lsernau/grh-lean/blob/72fead102f663d5c7a407964a13dc0cf9d977207/grh-compact/src/New/Data/Different.lean\#L321-L322}{{\leanfont Parameters.\allowbreak{}relDifferentGen\_ne\_zero}}}
\expandafter\newcommand\csname LN@difmul\endcsname{\href{https://github.com/lsernau/grh-lean/blob/72fead102f663d5c7a407964a13dc0cf9d977207/grh-compact/src/New/Data/Different.lean\#L324-L327}{{\leanfont Parameters.\allowbreak{}different\_eq\_span\_relDifferentGen\_mul}}}
\expandafter\newcommand\csname LN@condmem\endcsname{\href{https://github.com/lsernau/grh-lean/blob/72fead102f663d5c7a407964a13dc0cf9d977207/grh-compact/src/New/Data/Different.lean\#L330-L332}{{\leanfont Parameters.\allowbreak{}cycloConductor\_mem\_different}}}
\expandafter\newcommand\csname LN@unitreal\endcsname{\href{https://github.com/lsernau/grh-lean/blob/72fead102f663d5c7a407964a13dc0cf9d977207/grh-compact/src/New/Data/Different.lean\#L334-L342}{{\leanfont Parameters.\allowbreak{}exists\_unit\_mul\_relDifferentGen\_real}}}
\expandafter\newcommand\csname LN@vpi\endcsname{\href{https://github.com/lsernau/grh-lean/blob/72fead102f663d5c7a407964a13dc0cf9d977207/grh-compact/src/New/Data/Different.lean\#L346-L347}{{\leanfont Parameters.\allowbreak{}ϖ}}}
\expandafter\newcommand\csname LN@starvpi\endcsname{\href{https://github.com/lsernau/grh-lean/blob/72fead102f663d5c7a407964a13dc0cf9d977207/grh-compact/src/New/Data/Different.lean\#L350-L350}{{\leanfont Parameters.\allowbreak{}star\_ϖ}}}
\expandafter\newcommand\csname LN@difvpi\endcsname{\href{https://github.com/lsernau/grh-lean/blob/72fead102f663d5c7a407964a13dc0cf9d977207/grh-compact/src/New/Data/Different.lean\#L353-L354}{{\leanfont Parameters.\allowbreak{}different\_eq\_span\_varpi}}}
\expandafter\newcommand\csname LN@prodvpi\endcsname{\href{https://github.com/lsernau/grh-lean/blob/72fead102f663d5c7a407964a13dc0cf9d977207/grh-compact/src/New/Data/Different.lean\#L358-L362}{{\leanfont Parameters.\allowbreak{}prod\_varpi\_eq\_rootDiscr\_pow}}}
\expandafter\newcommand\csname LQ@prodvpi\endcsname{\leanquoteline{0}{lemma prod\_varpi\_eq\_rootDiscr\_pow : ∏ v : P.Φ s, v (P.ϖ s).val =}\leanquoteline{2}{rootDiscr (P.E s) \textasciicircum{} numInfPlaces (P.E s)}}
\expandafter\newcommand\csname LN@absvpi\endcsname{\href{https://github.com/lsernau/grh-lean/blob/72fead102f663d5c7a407964a13dc0cf9d977207/grh-compact/src/New/Data/Different.lean\#L364-L366}{{\leanfont Parameters.\allowbreak{}infinitePlace\_apply\_varpi}}}
\expandafter\newcommand\csname LN@galv\endcsname{\href{https://github.com/lsernau/grh-lean/blob/72fead102f663d5c7a407964a13dc0cf9d977207/grh-compact/src/New/Data/WeilDataOfE.lean\#L24-L24}{{\leanfont Parameters.\allowbreak{}galOfPlaceInt}}}
\expandafter\newcommand\csname LN@coegal\endcsname{\href{https://github.com/lsernau/grh-lean/blob/72fead102f663d5c7a407964a13dc0cf9d977207/grh-compact/src/New/Data/WeilDataOfE.lean\#L26-L26}{{\leanfont Parameters.\allowbreak{}coe\_galOfPlaceInt}}}
\expandafter\newcommand\csname LN@memsmul\endcsname{\href{https://github.com/lsernau/grh-lean/blob/72fead102f663d5c7a407964a13dc0cf9d977207/grh-compact/src/New/Data/WeilDataOfE.lean\#L32-L32}{{\leanfont Parameters.\allowbreak{}mem\_smul}}}
\expandafter\newcommand\csname LN@Frb\endcsname{\href{https://github.com/lsernau/grh-lean/blob/72fead102f663d5c7a407964a13dc0cf9d977207/grh-compact/src/New/Data/WeilDataOfE.lean\#L36-L36}{{\leanfont Parameters.\allowbreak{}frobElem}}}
\expandafter\newcommand\csname LN@copNz\endcsname{\href{https://github.com/lsernau/grh-lean/blob/72fead102f663d5c7a407964a13dc0cf9d977207/grh-compact/src/New/Data/WeilDataOfE.lean\#L39-L40}{{\leanfont Parameters.\allowbreak{}isCoprime\_N₀\_of\_not\_dvd}}}
\expandafter\newcommand\csname LN@frobext\endcsname{\href{https://github.com/lsernau/grh-lean/blob/72fead102f663d5c7a407964a13dc0cf9d977207/grh-compact/src/New/Data/WeilDataOfE.lean\#L46-L56}{{\leanfont Parameters.\allowbreak{}frob\_ext}}}
\expandafter\newcommand\csname LN@frobpow\endcsname{\href{https://github.com/lsernau/grh-lean/blob/72fead102f663d5c7a407964a13dc0cf9d977207/grh-compact/src/New/Data/WeilDataOfE.lean\#L59-L69}{{\leanfont Parameters.\allowbreak{}frob\_iff\_eq\_frobElem\_pow}}}
\expandafter\newcommand\csname LN@decomp\endcsname{\href{https://github.com/lsernau/grh-lean/blob/72fead102f663d5c7a407964a13dc0cf9d977207/grh-compact/src/New/Data/WeilDataOfE.lean\#L77-L84}{{\leanfont Parameters.\allowbreak{}smul\_eq\_iff\_exists\_frobElem\_pow}}}
\expandafter\newcommand\csname LQ@decomp\endcsname{\leanquoteline{0}{lemma smul\_eq\_iff\_exists\_frobElem\_pow : v • w = comap c c.surjective w ↔ ∃ n, (v, e) =}\leanquoteline{2}{P.frobElem s p \textasciicircum{} n}}
\expandafter\newcommand\csname LN@qpow\endcsname{\href{https://github.com/lsernau/grh-lean/blob/72fead102f663d5c7a407964a13dc0cf9d977207/grh-compact/src/New/Data/WeilDataOfE.lean\#L88-L91}{{\leanfont Parameters.\allowbreak{}q\_eq\_pow}}}
\expandafter\newcommand\csname LN@notsq\endcsname{\href{https://github.com/lsernau/grh-lean/blob/72fead102f663d5c7a407964a13dc0cf9d977207/grh-compact/src/New/Data/WeilDataOfE.lean\#L94-L95}{{\leanfont Parameters.\allowbreak{}natCast\_notMem\_sq}}}
\expandafter\newcommand\csname LN@valp\endcsname{\href{https://github.com/lsernau/grh-lean/blob/72fead102f663d5c7a407964a13dc0cf9d977207/grh-compact/src/New/Data/WeilDataOfE.lean\#L100-L101}{{\leanfont Parameters.\allowbreak{}intValuation\_p}}}
\expandafter\newcommand\csname LN@transit\endcsname{\href{https://github.com/lsernau/grh-lean/blob/72fead102f663d5c7a407964a13dc0cf9d977207/grh-compact/src/New/Data/WeilDataOfE.lean\#L105-L110}{{\leanfont Parameters.\allowbreak{}exists\_smul\_eq\_or\_star}}}
\expandafter\newcommand\csname LN@inertnd\endcsname{\href{https://github.com/lsernau/grh-lean/blob/72fead102f663d5c7a407964a13dc0cf9d977207/grh-compact/src/New/Data/WeilDataOfE.lean\#L112-L113}{{\leanfont Parameters.\allowbreak{}not\_dvd\_N₁\_mul\_d\_of\_inert}}}
\expandafter\newcommand\csname LN@inertfix\endcsname{\href{https://github.com/lsernau/grh-lean/blob/72fead102f663d5c7a407964a13dc0cf9d977207/grh-compact/src/New/Data/WeilDataOfE.lean\#L119-L120}{{\leanfont Parameters.\allowbreak{}smul\_eq\_iff\_inert}}}
\expandafter\newcommand\csname LN@inertstar\endcsname{\href{https://github.com/lsernau/grh-lean/blob/72fead102f663d5c7a407964a13dc0cf9d977207/grh-compact/src/New/Data/WeilDataOfE.lean\#L124-L125}{{\leanfont Parameters.\allowbreak{}smul\_eq\_star\_iff\_inert}}}
\expandafter\newcommand\csname LN@inertq\endcsname{\href{https://github.com/lsernau/grh-lean/blob/72fead102f663d5c7a407964a13dc0cf9d977207/grh-compact/src/New/Data/WeilDataOfE.lean\#L129-L130}{{\leanfont Parameters.\allowbreak{}q\_eq\_inert}}}
\expandafter\newcommand\csname LN@splitchi\endcsname{\href{https://github.com/lsernau/grh-lean/blob/72fead102f663d5c7a407964a13dc0cf9d977207/grh-compact/src/New/Data/WeilDataOfE.lean\#L137-L137}{{\leanfont Parameters.\allowbreak{}chiK\_eq\_one\_of\_modEq\_one}}}
\expandafter\newcommand\csname LN@splitnd\endcsname{\href{https://github.com/lsernau/grh-lean/blob/72fead102f663d5c7a407964a13dc0cf9d977207/grh-compact/src/New/Data/WeilDataOfE.lean\#L139-L140}{{\leanfont Parameters.\allowbreak{}not\_dvd\_N₁\_mul\_d\_of\_modEq\_one}}}
\expandafter\newcommand\csname LN@splitfr\endcsname{\href{https://github.com/lsernau/grh-lean/blob/72fead102f663d5c7a407964a13dc0cf9d977207/grh-compact/src/New/Data/WeilDataOfE.lean\#L144-L145}{{\leanfont Parameters.\allowbreak{}frobElem\_eq\_one}}}
\expandafter\newcommand\csname LN@splitne\endcsname{\href{https://github.com/lsernau/grh-lean/blob/72fead102f663d5c7a407964a13dc0cf9d977207/grh-compact/src/New/Data/WeilDataOfE.lean\#L151-L152}{{\leanfont Parameters.\allowbreak{}smul\_ne\_star}}}
\expandafter\newcommand\csname LN@splitfix\endcsname{\href{https://github.com/lsernau/grh-lean/blob/72fead102f663d5c7a407964a13dc0cf9d977207/grh-compact/src/New/Data/WeilDataOfE.lean\#L154-L155}{{\leanfont Parameters.\allowbreak{}smul\_eq\_self\_iff\_eq\_one}}}
\expandafter\newcommand\csname LN@splitq\endcsname{\href{https://github.com/lsernau/grh-lean/blob/72fead102f663d5c7a407964a13dc0cf9d977207/grh-compact/src/New/Data/WeilDataOfE.lean\#L158-L158}{{\leanfont Parameters.\allowbreak{}q\_eq\_of\_modEq\_one}}}
\expandafter\newcommand\csname LN@starmem\endcsname{\href{https://github.com/lsernau/grh-lean/blob/72fead102f663d5c7a407964a13dc0cf9d977207/grh-compact/src/New/Data/WeilDataOfE.lean\#L164-L166}{{\leanfont Parameters.\allowbreak{}star\_mem\_iff}}}
\expandafter\newcommand\csname LN@valnorm\endcsname{\href{https://github.com/lsernau/grh-lean/blob/72fead102f663d5c7a407964a13dc0cf9d977207/grh-compact/src/New/Data/WeilDataOfE.lean\#L170-L172}{{\leanfont Parameters.\allowbreak{}intValuation\_eq\_of\_mul\_star}}}
\expandafter\newcommand\csname LN@bij\endcsname{\href{https://github.com/lsernau/grh-lean/blob/72fead102f663d5c7a407964a13dc0cf9d977207/grh-compact/src/New/Data/WeilDataOfE.lean\#L178-L183}{{\leanfont Parameters.\allowbreak{}nonempty\_infinitePlace\_equiv\_finPlace}}}
\expandafter\newcommand\csname LN@valmap\endcsname{\href{https://github.com/lsernau/grh-lean/blob/72fead102f663d5c7a407964a13dc0cf9d977207/grh-compact/src/New/Data/WeilDataOfE.lean\#L192-L196}{{\leanfont Parameters.\allowbreak{}valuation\_map\_relDifferentGen}}}
\expandafter\newcommand\csname LN@alP\endcsname{\href{https://github.com/lsernau/grh-lean/blob/72fead102f663d5c7a407964a13dc0cf9d977207/grh-compact/src/New/Data/WeilDataOfE.lean\#L200-L202}{{\leanfont Parameters.\allowbreak{}α}}}
\expandafter\newcommand\csname LN@alinv\endcsname{\href{https://github.com/lsernau/grh-lean/blob/72fead102f663d5c7a407964a13dc0cf9d977207/grh-compact/src/New/Data/WeilDataOfE.lean\#L207-L207}{{\leanfont Parameters.\allowbreak{}inv\_α\_mem}}}
\expandafter\newcommand\csname LN@valpha\endcsname{\href{https://github.com/lsernau/grh-lean/blob/72fead102f663d5c7a407964a13dc0cf9d977207/grh-compact/src/New/Data/WeilDataOfE.lean\#L211-L217}{{\leanfont Parameters.\allowbreak{}v\_alpha\_mul\_relDifferentGen}}}
\expandafter\newcommand\csname LN@Nal\endcsname{\href{https://github.com/lsernau/grh-lean/blob/72fead102f663d5c7a407964a13dc0cf9d977207/grh-compact/src/New/Data/WeilDataOfE.lean\#L219-L220}{{\leanfont Parameters.\allowbreak{}N₁\_mul\_α\_mem}}}
\expandafter\newcommand\csname LN@weil\endcsname{\href{https://github.com/lsernau/grh-lean/blob/72fead102f663d5c7a407964a13dc0cf9d977207/grh-compact/src/New/Data/WeilDataOfE.lean\#L225-L243}{{\leanfont Parameters.\allowbreak{}instWeilDataE}}}
\newcommand{\ConcordanceRows}{\crow{par}{\href{https://github.com/lsernau/grh-lean/blob/72fead102f663d5c7a407964a13dc0cf9d977207/grh-compact/src/New/Data/Different.lean\#L23-L28}{{\leanfont Parameters}}}{Different.lean}{23--28}
\crow{twd}{\href{https://github.com/lsernau/grh-lean/blob/72fead102f663d5c7a407964a13dc0cf9d977207/grh-compact/src/New/Data/Different.lean\#L41-L41}{{\leanfont towerDeg}}}{Different.lean}{41--41}
\crow{Nn}{\href{https://github.com/lsernau/grh-lean/blob/72fead102f663d5c7a407964a13dc0cf9d977207/grh-compact/src/New/Data/Different.lean\#L44-L44}{{\leanfont Parameters.\allowbreak{}N₁}}}{Different.lean}{44--44}
\crow{KP}{\href{https://github.com/lsernau/grh-lean/blob/72fead102f663d5c7a407964a13dc0cf9d977207/grh-compact/src/New/Data/Different.lean\#L64-L64}{{\leanfont Parameters.\allowbreak{}K}}}{Different.lean}{64--64}
\crow{Mom}{\href{https://github.com/lsernau/grh-lean/blob/72fead102f663d5c7a407964a13dc0cf9d977207/grh-compact/src/New/Data/Different.lean\#L68-L68}{{\leanfont Parameters.\allowbreak{}Mω}}}{Different.lean}{68--68}
\crow{heeg}{\href{https://github.com/lsernau/grh-lean/blob/72fead102f663d5c7a407964a13dc0cf9d977207/grh-compact/src/New/Data/Different.lean\#L36-L37}{{\leanfont Parameters.\allowbreak{}instHeegD}}}{Different.lean}{36--37}
\crow{ltd}{\href{https://github.com/lsernau/grh-lean/blob/72fead102f663d5c7a407964a13dc0cf9d977207/grh-compact/src/New/Data/Different.lean\#L52-L52}{{\leanfont Parameters.\allowbreak{}lt\_d}}}{Different.lean}{52--52}
\crow{dvdnorm}{\href{https://github.com/lsernau/grh-lean/blob/72fead102f663d5c7a407964a13dc0cf9d977207/grh-compact/src/New/Data/Different.lean\#L47-L50}{{\leanfont Parameters.\allowbreak{}dvd\_norm\_c₀}}}{Different.lean}{47--50}
\crow{chione}{\href{https://github.com/lsernau/grh-lean/blob/72fead102f663d5c7a407964a13dc0cf9d977207/grh-compact/src/New/Data/Different.lean\#L55-L58}{{\leanfont Parameters.\allowbreak{}chiK\_eq\_one\_of\_dvd}}}{Different.lean}{55--58}
\crow{thirteen}{\href{https://github.com/lsernau/grh-lean/blob/72fead102f663d5c7a407964a13dc0cf9d977207/grh-compact/src/New/Data/Different.lean\#L60-L61}{{\leanfont Parameters.\allowbreak{}thirteen\_le\_of\_chiK\_ne\_one}}}{Different.lean}{60--61}
\crow{NdM}{\href{https://github.com/lsernau/grh-lean/blob/72fead102f663d5c7a407964a13dc0cf9d977207/grh-compact/src/New/Data/Different.lean\#L71-L72}{{\leanfont Parameters.\allowbreak{}N₁\_mul\_d\_dvd\_M\_omega}}}{Different.lean}{71--72}
\crow{coprime}{\href{https://github.com/lsernau/grh-lean/blob/72fead102f663d5c7a407964a13dc0cf9d977207/grh-compact/src/New/Data/Different.lean\#L74-L75}{{\leanfont Parameters.\allowbreak{}coprime\_N₁\_d}}}{Different.lean}{74--75}
\crow{twopow}{\href{https://github.com/lsernau/grh-lean/blob/72fead102f663d5c7a407964a13dc0cf9d977207/grh-compact/src/New/Data/Different.lean\#L193-L194}{{\leanfont Parameters.\allowbreak{}two\_pow\_dvd\_N₁}}}{Different.lean}{193--194}
\crow{powd}{\href{https://github.com/lsernau/grh-lean/blob/72fead102f663d5c7a407964a13dc0cf9d977207/grh-compact/src/New/Data/Different.lean\#L212-L215}{{\leanfont Parameters.\allowbreak{}pow\_d\_ne\_neg\_one}}}{Different.lean}{212--215}
\crow{EP}{\href{https://github.com/lsernau/grh-lean/blob/72fead102f663d5c7a407964a13dc0cf9d977207/grh-compact/src/New/Data/Different.lean\#L79-L79}{{\leanfont Parameters.\allowbreak{}E}}}{Different.lean}{79--79}
\crow{PlP}{\href{https://github.com/lsernau/grh-lean/blob/72fead102f663d5c7a407964a13dc0cf9d977207/grh-compact/src/New/Data/Different.lean\#L83-L83}{{\leanfont Parameters.\allowbreak{}Φ}}}{Different.lean}{83--83}
\crow{rtu}{\href{https://github.com/lsernau/grh-lean/blob/72fead102f663d5c7a407964a13dc0cf9d977207/grh-compact/src/New/Data/Different.lean\#L142-L142}{{\leanfont Parameters.\allowbreak{}ζ}}}{Different.lean}{142--142}
\crow{OmP}{\href{https://github.com/lsernau/grh-lean/blob/72fead102f663d5c7a407964a13dc0cf9d977207/grh-compact/src/New/Data/Different.lean\#L87-L87}{{\leanfont Parameters.\allowbreak{}bigOmega}}}{Different.lean}{87--87}
\crow{primrt}{\href{https://github.com/lsernau/grh-lean/blob/72fead102f663d5c7a407964a13dc0cf9d977207/grh-compact/src/New/Data/Different.lean\#L180-L181}{{\leanfont Parameters.\allowbreak{}isPrimitiveRoot\_zeta}}}{Different.lean}{180--181}
\crow{irr}{\href{https://github.com/lsernau/grh-lean/blob/72fead102f663d5c7a407964a13dc0cf9d977207/grh-compact/src/New/Data/Different.lean\#L92-L100}{{\leanfont Parameters.\allowbreak{}irreducible\_cyclotomic}}}{Different.lean}{92--100}
\crow{tocyc}{\href{https://github.com/lsernau/grh-lean/blob/72fead102f663d5c7a407964a13dc0cf9d977207/grh-compact/src/New/Data/Different.lean\#L104-L108}{{\leanfont Parameters.\allowbreak{}toCycloField}}}{Different.lean}{104--108}
\crow{embeq}{\href{https://github.com/lsernau/grh-lean/blob/72fead102f663d5c7a407964a13dc0cf9d977207/grh-compact/src/New/Data/Different.lean\#L116-L117}{{\leanfont Parameters.\allowbreak{}algHom\_eq\_of\_mk\_eq}}}{Different.lean}{116--117}
\crow{galP}{\href{https://github.com/lsernau/grh-lean/blob/72fead102f663d5c7a407964a13dc0cf9d977207/grh-compact/src/New/Data/Different.lean\#L127-L127}{{\leanfont Parameters.\allowbreak{}galOfPlace}}}{Different.lean}{127--127}
\crow{cycP}{\href{https://github.com/lsernau/grh-lean/blob/72fead102f663d5c7a407964a13dc0cf9d977207/grh-compact/src/New/Data/Different.lean\#L131-L131}{{\leanfont Parameters.\allowbreak{}placeMulEquiv}}}{Different.lean}{131--131}
\crow{pin}{\href{https://github.com/lsernau/grh-lean/blob/72fead102f663d5c7a407964a13dc0cf9d977207/grh-compact/src/New/Data/Different.lean\#L135-L135}{{\leanfont Parameters.\allowbreak{}placeOfInt}}}{Different.lean}{135--135}
\crow{pcl}{\href{https://github.com/lsernau/grh-lean/blob/72fead102f663d5c7a407964a13dc0cf9d977207/grh-compact/src/New/Data/Different.lean\#L139-L139}{{\leanfont Parameters.\allowbreak{}placeClass}}}{Different.lean}{139--139}
\crow{phialg}{\href{https://github.com/lsernau/grh-lean/blob/72fead102f663d5c7a407964a13dc0cf9d977207/grh-compact/src/New/Data/Different.lean\#L148-L148}{{\leanfont Parameters.\allowbreak{}phi\_algebraMap}}}{Different.lean}{148--148}
\crow{phimul}{\href{https://github.com/lsernau/grh-lean/blob/72fead102f663d5c7a407964a13dc0cf9d977207/grh-compact/src/New/Data/Different.lean\#L183-L183}{{\leanfont Parameters.\allowbreak{}phi\_mul}}}{Different.lean}{183--183}
\crow{starrt}{\href{https://github.com/lsernau/grh-lean/blob/72fead102f663d5c7a407964a13dc0cf9d977207/grh-compact/src/New/Data/Different.lean\#L185-L187}{{\leanfont Parameters.\allowbreak{}star\_ζ}}}{Different.lean}{185--187}
\crow{galrt}{\href{https://github.com/lsernau/grh-lean/blob/72fead102f663d5c7a407964a13dc0cf9d977207/grh-compact/src/New/Data/Different.lean\#L189-L190}{{\leanfont Parameters.\allowbreak{}galOfPlace\_zeta}}}{Different.lean}{189--190}
\crow{pinpow}{\href{https://github.com/lsernau/grh-lean/blob/72fead102f663d5c7a407964a13dc0cf9d977207/grh-compact/src/New/Data/Different.lean\#L196-L197}{{\leanfont Parameters.\allowbreak{}placeOfInt\_pow}}}{Different.lean}{196--197}
\crow{clpin}{\href{https://github.com/lsernau/grh-lean/blob/72fead102f663d5c7a407964a13dc0cf9d977207/grh-compact/src/New/Data/Different.lean\#L199-L202}{{\leanfont Parameters.\allowbreak{}placeClass\_placeOfInt}}}{Different.lean}{199--202}
\crow{galpin}{\href{https://github.com/lsernau/grh-lean/blob/72fead102f663d5c7a407964a13dc0cf9d977207/grh-compact/src/New/Data/Different.lean\#L206-L209}{{\leanfont Parameters.\allowbreak{}galOfPlace\_placeOfInt\_zeta}}}{Different.lean}{206--209}
\crow{clneg}{\href{https://github.com/lsernau/grh-lean/blob/72fead102f663d5c7a407964a13dc0cf9d977207/grh-compact/src/New/Data/Different.lean\#L217-L220}{{\leanfont Parameters.\allowbreak{}placeClass\_placeOfInt\_eq\_neg\_one}}}{Different.lean}{217--220}
\crow{expandpow}{\href{https://github.com/lsernau/grh-lean/blob/72fead102f663d5c7a407964a13dc0cf9d977207/grh-compact/src/New/Data/Different.lean\#L156-L158}{{\leanfont Parameters.\allowbreak{}expand\_pow\_cyclotomic}}}{Different.lean}{156--158}
\crow{expandtw}{\href{https://github.com/lsernau/grh-lean/blob/72fead102f663d5c7a407964a13dc0cf9d977207/grh-compact/src/New/Data/Different.lean\#L162-L164}{{\leanfont Parameters.\allowbreak{}expand\_towerDeg\_cyclotomic}}}{Different.lean}{162--164}
\crow{cardPl}{\href{https://github.com/lsernau/grh-lean/blob/72fead102f663d5c7a407964a13dc0cf9d977207/grh-compact/src/New/Data/Different.lean\#L168-L170}{{\leanfont Parameters.\allowbreak{}card\_infinitePlace\_E}}}{Different.lean}{168--170}
\crow{cardcl}{\href{https://github.com/lsernau/grh-lean/blob/72fead102f663d5c7a407964a13dc0cf9d977207/grh-compact/src/New/Data/Different.lean\#L174-L176}{{\leanfont Parameters.\allowbreak{}card\_placeClass\_eq}}}{Different.lean}{174--176}
\crow{prm}{\href{https://github.com/lsernau/grh-lean/blob/72fead102f663d5c7a407964a13dc0cf9d977207/grh-compact/src/New/Data/Different.lean\#L224-L224}{{\leanfont Parameters.\allowbreak{}placePerm}}}{Different.lean}{224--224}
\crow{prmspec}{\href{https://github.com/lsernau/grh-lean/blob/72fead102f663d5c7a407964a13dc0cf9d977207/grh-compact/src/New/Data/Different.lean\#L229-L231}{{\leanfont Parameters.\allowbreak{}placePerm\_spec}}}{Different.lean}{229--231}
\crow{roots}{\href{https://github.com/lsernau/grh-lean/blob/72fead102f663d5c7a407964a13dc0cf9d977207/grh-compact/src/New/Data/Different.lean\#L239-L247}{{\leanfont Parameters.\allowbreak{}exists\_zeta\_pow\_eq}}}{Different.lean}{239--247}
\crow{prodphi}{\href{https://github.com/lsernau/grh-lean/blob/72fead102f663d5c7a407964a13dc0cf9d977207/grh-compact/src/New/Data/Different.lean\#L249-L250}{{\leanfont Parameters.\allowbreak{}prod\_phi\_eq\_norm}}}{Different.lean}{249--250}
\crow{rti}{\href{https://github.com/lsernau/grh-lean/blob/72fead102f663d5c7a407964a13dc0cf9d977207/grh-compact/src/New/Data/Different.lean\#L260-L260}{{\leanfont Parameters.\allowbreak{}zetaInt}}}{Different.lean}{260--260}
\crow{czE}{\href{https://github.com/lsernau/grh-lean/blob/72fead102f663d5c7a407964a13dc0cf9d977207/grh-compact/src/New/Data/Different.lean\#L268-L268}{{\leanfont Parameters.\allowbreak{}c₀InE}}}{Different.lean}{268--268}
\crow{sqd}{\href{https://github.com/lsernau/grh-lean/blob/72fead102f663d5c7a407964a13dc0cf9d977207/grh-compact/src/New/Data/Different.lean\#L271-L271}{{\leanfont Parameters.\allowbreak{}sqrtNegDInE}}}{Different.lean}{271--271}
\crow{dgen}{\href{https://github.com/lsernau/grh-lean/blob/72fead102f663d5c7a407964a13dc0cf9d977207/grh-compact/src/New/Data/Different.lean\#L317-L317}{{\leanfont Parameters.\allowbreak{}relDifferentGen}}}{Different.lean}{317--317}
\crow{coert}{\href{https://github.com/lsernau/grh-lean/blob/72fead102f663d5c7a407964a13dc0cf9d977207/grh-compact/src/New/Data/Different.lean\#L262-L262}{{\leanfont Parameters.\allowbreak{}coe\_zetaInt}}}{Different.lean}{262--262}
\crow{primrti}{\href{https://github.com/lsernau/grh-lean/blob/72fead102f663d5c7a407964a13dc0cf9d977207/grh-compact/src/New/Data/Different.lean\#L264-L265}{{\leanfont Parameters.\allowbreak{}isPrimitiveRoot\_zetaInt}}}{Different.lean}{264--265}
\crow{coesq}{\href{https://github.com/lsernau/grh-lean/blob/72fead102f663d5c7a407964a13dc0cf9d977207/grh-compact/src/New/Data/Different.lean\#L276-L277}{{\leanfont Parameters.\allowbreak{}coe\_sqrtNegDInE}}}{Different.lean}{276--277}
\crow{starsq}{\href{https://github.com/lsernau/grh-lean/blob/72fead102f663d5c7a407964a13dc0cf9d977207/grh-compact/src/New/Data/Different.lean\#L279-L280}{{\leanfont Parameters.\allowbreak{}star\_sqrtNegDInE}}}{Different.lean}{279--280}
\crow{starrtpow}{\href{https://github.com/lsernau/grh-lean/blob/72fead102f663d5c7a407964a13dc0cf9d977207/grh-compact/src/New/Data/Different.lean\#L282-L283}{{\leanfont Parameters.\allowbreak{}star\_zetaInt\_pow\_mul}}}{Different.lean}{282--283}
\crow{sqdvd}{\href{https://github.com/lsernau/grh-lean/blob/72fead102f663d5c7a407964a13dc0cf9d977207/grh-compact/src/New/Data/Different.lean\#L285-L286}{{\leanfont Parameters.\allowbreak{}sqrtNegDInE\_dvd}}}{Different.lean}{285--286}
\crow{czmul}{\href{https://github.com/lsernau/grh-lean/blob/72fead102f663d5c7a407964a13dc0cf9d977207/grh-compact/src/New/Data/Different.lean\#L288-L289}{{\leanfont Parameters.\allowbreak{}c₀InE\_mul\_one\_sub}}}{Different.lean}{288--289}
\crow{starcz}{\href{https://github.com/lsernau/grh-lean/blob/72fead102f663d5c7a407964a13dc0cf9d977207/grh-compact/src/New/Data/Different.lean\#L291-L292}{{\leanfont Parameters.\allowbreak{}star\_c₀InE}}}{Different.lean}{291--292}
\crow{relspec}{\href{https://github.com/lsernau/grh-lean/blob/72fead102f663d5c7a407964a13dc0cf9d977207/grh-compact/src/New/Data/Different.lean\#L299-L313}{{\leanfont Parameters.\allowbreak{}relDifferent\_spec}}}{Different.lean}{299--313}
\crow{dgendvd}{\href{https://github.com/lsernau/grh-lean/blob/72fead102f663d5c7a407964a13dc0cf9d977207/grh-compact/src/New/Data/Different.lean\#L319-L319}{{\leanfont Parameters.\allowbreak{}relDifferentGen\_dvd\_N₁}}}{Different.lean}{319--319}
\crow{dgenne}{\href{https://github.com/lsernau/grh-lean/blob/72fead102f663d5c7a407964a13dc0cf9d977207/grh-compact/src/New/Data/Different.lean\#L321-L322}{{\leanfont Parameters.\allowbreak{}relDifferentGen\_ne\_zero}}}{Different.lean}{321--322}
\crow{difmul}{\href{https://github.com/lsernau/grh-lean/blob/72fead102f663d5c7a407964a13dc0cf9d977207/grh-compact/src/New/Data/Different.lean\#L324-L327}{{\leanfont Parameters.\allowbreak{}different\_eq\_span\_relDifferentGen\_mul}}}{Different.lean}{324--327}
\crow{condmem}{\href{https://github.com/lsernau/grh-lean/blob/72fead102f663d5c7a407964a13dc0cf9d977207/grh-compact/src/New/Data/Different.lean\#L330-L332}{{\leanfont Parameters.\allowbreak{}cycloConductor\_mem\_different}}}{Different.lean}{330--332}
\crow{unitreal}{\href{https://github.com/lsernau/grh-lean/blob/72fead102f663d5c7a407964a13dc0cf9d977207/grh-compact/src/New/Data/Different.lean\#L334-L342}{{\leanfont Parameters.\allowbreak{}exists\_unit\_mul\_relDifferentGen\_real}}}{Different.lean}{334--342}
\crow{vpi}{\href{https://github.com/lsernau/grh-lean/blob/72fead102f663d5c7a407964a13dc0cf9d977207/grh-compact/src/New/Data/Different.lean\#L346-L347}{{\leanfont Parameters.\allowbreak{}ϖ}}}{Different.lean}{346--347}
\crow{starvpi}{\href{https://github.com/lsernau/grh-lean/blob/72fead102f663d5c7a407964a13dc0cf9d977207/grh-compact/src/New/Data/Different.lean\#L350-L350}{{\leanfont Parameters.\allowbreak{}star\_ϖ}}}{Different.lean}{350--350}
\crow{difvpi}{\href{https://github.com/lsernau/grh-lean/blob/72fead102f663d5c7a407964a13dc0cf9d977207/grh-compact/src/New/Data/Different.lean\#L353-L354}{{\leanfont Parameters.\allowbreak{}different\_eq\_span\_varpi}}}{Different.lean}{353--354}
\crow{prodvpi}{\href{https://github.com/lsernau/grh-lean/blob/72fead102f663d5c7a407964a13dc0cf9d977207/grh-compact/src/New/Data/Different.lean\#L358-L362}{{\leanfont Parameters.\allowbreak{}prod\_varpi\_eq\_rootDiscr\_pow}}}{Different.lean}{358--362}
\crow{absvpi}{\href{https://github.com/lsernau/grh-lean/blob/72fead102f663d5c7a407964a13dc0cf9d977207/grh-compact/src/New/Data/Different.lean\#L364-L366}{{\leanfont Parameters.\allowbreak{}infinitePlace\_apply\_varpi}}}{Different.lean}{364--366}
\crow{galv}{\href{https://github.com/lsernau/grh-lean/blob/72fead102f663d5c7a407964a13dc0cf9d977207/grh-compact/src/New/Data/WeilDataOfE.lean\#L24-L24}{{\leanfont Parameters.\allowbreak{}galOfPlaceInt}}}{WeilDataOfE.lean}{24--24}
\crow{coegal}{\href{https://github.com/lsernau/grh-lean/blob/72fead102f663d5c7a407964a13dc0cf9d977207/grh-compact/src/New/Data/WeilDataOfE.lean\#L26-L26}{{\leanfont Parameters.\allowbreak{}coe\_galOfPlaceInt}}}{WeilDataOfE.lean}{26--26}
\crow{memsmul}{\href{https://github.com/lsernau/grh-lean/blob/72fead102f663d5c7a407964a13dc0cf9d977207/grh-compact/src/New/Data/WeilDataOfE.lean\#L32-L32}{{\leanfont Parameters.\allowbreak{}mem\_smul}}}{WeilDataOfE.lean}{32--32}
\crow{Frb}{\href{https://github.com/lsernau/grh-lean/blob/72fead102f663d5c7a407964a13dc0cf9d977207/grh-compact/src/New/Data/WeilDataOfE.lean\#L36-L36}{{\leanfont Parameters.\allowbreak{}frobElem}}}{WeilDataOfE.lean}{36--36}
\crow{copNz}{\href{https://github.com/lsernau/grh-lean/blob/72fead102f663d5c7a407964a13dc0cf9d977207/grh-compact/src/New/Data/WeilDataOfE.lean\#L39-L40}{{\leanfont Parameters.\allowbreak{}isCoprime\_N₀\_of\_not\_dvd}}}{WeilDataOfE.lean}{39--40}
\crow{frobext}{\href{https://github.com/lsernau/grh-lean/blob/72fead102f663d5c7a407964a13dc0cf9d977207/grh-compact/src/New/Data/WeilDataOfE.lean\#L46-L56}{{\leanfont Parameters.\allowbreak{}frob\_ext}}}{WeilDataOfE.lean}{46--56}
\crow{frobpow}{\href{https://github.com/lsernau/grh-lean/blob/72fead102f663d5c7a407964a13dc0cf9d977207/grh-compact/src/New/Data/WeilDataOfE.lean\#L59-L69}{{\leanfont Parameters.\allowbreak{}frob\_iff\_eq\_frobElem\_pow}}}{WeilDataOfE.lean}{59--69}
\crow{decomp}{\href{https://github.com/lsernau/grh-lean/blob/72fead102f663d5c7a407964a13dc0cf9d977207/grh-compact/src/New/Data/WeilDataOfE.lean\#L77-L84}{{\leanfont Parameters.\allowbreak{}smul\_eq\_iff\_exists\_frobElem\_pow}}}{WeilDataOfE.lean}{77--84}
\crow{qpow}{\href{https://github.com/lsernau/grh-lean/blob/72fead102f663d5c7a407964a13dc0cf9d977207/grh-compact/src/New/Data/WeilDataOfE.lean\#L88-L91}{{\leanfont Parameters.\allowbreak{}q\_eq\_pow}}}{WeilDataOfE.lean}{88--91}
\crow{notsq}{\href{https://github.com/lsernau/grh-lean/blob/72fead102f663d5c7a407964a13dc0cf9d977207/grh-compact/src/New/Data/WeilDataOfE.lean\#L94-L95}{{\leanfont Parameters.\allowbreak{}natCast\_notMem\_sq}}}{WeilDataOfE.lean}{94--95}
\crow{valp}{\href{https://github.com/lsernau/grh-lean/blob/72fead102f663d5c7a407964a13dc0cf9d977207/grh-compact/src/New/Data/WeilDataOfE.lean\#L100-L101}{{\leanfont Parameters.\allowbreak{}intValuation\_p}}}{WeilDataOfE.lean}{100--101}
\crow{transit}{\href{https://github.com/lsernau/grh-lean/blob/72fead102f663d5c7a407964a13dc0cf9d977207/grh-compact/src/New/Data/WeilDataOfE.lean\#L105-L110}{{\leanfont Parameters.\allowbreak{}exists\_smul\_eq\_or\_star}}}{WeilDataOfE.lean}{105--110}
\crow{inertnd}{\href{https://github.com/lsernau/grh-lean/blob/72fead102f663d5c7a407964a13dc0cf9d977207/grh-compact/src/New/Data/WeilDataOfE.lean\#L112-L113}{{\leanfont Parameters.\allowbreak{}not\_dvd\_N₁\_mul\_d\_of\_inert}}}{WeilDataOfE.lean}{112--113}
\crow{inertfix}{\href{https://github.com/lsernau/grh-lean/blob/72fead102f663d5c7a407964a13dc0cf9d977207/grh-compact/src/New/Data/WeilDataOfE.lean\#L119-L120}{{\leanfont Parameters.\allowbreak{}smul\_eq\_iff\_inert}}}{WeilDataOfE.lean}{119--120}
\crow{inertstar}{\href{https://github.com/lsernau/grh-lean/blob/72fead102f663d5c7a407964a13dc0cf9d977207/grh-compact/src/New/Data/WeilDataOfE.lean\#L124-L125}{{\leanfont Parameters.\allowbreak{}smul\_eq\_star\_iff\_inert}}}{WeilDataOfE.lean}{124--125}
\crow{inertq}{\href{https://github.com/lsernau/grh-lean/blob/72fead102f663d5c7a407964a13dc0cf9d977207/grh-compact/src/New/Data/WeilDataOfE.lean\#L129-L130}{{\leanfont Parameters.\allowbreak{}q\_eq\_inert}}}{WeilDataOfE.lean}{129--130}
\crow{splitchi}{\href{https://github.com/lsernau/grh-lean/blob/72fead102f663d5c7a407964a13dc0cf9d977207/grh-compact/src/New/Data/WeilDataOfE.lean\#L137-L137}{{\leanfont Parameters.\allowbreak{}chiK\_eq\_one\_of\_modEq\_one}}}{WeilDataOfE.lean}{137--137}
\crow{splitnd}{\href{https://github.com/lsernau/grh-lean/blob/72fead102f663d5c7a407964a13dc0cf9d977207/grh-compact/src/New/Data/WeilDataOfE.lean\#L139-L140}{{\leanfont Parameters.\allowbreak{}not\_dvd\_N₁\_mul\_d\_of\_modEq\_one}}}{WeilDataOfE.lean}{139--140}
\crow{splitfr}{\href{https://github.com/lsernau/grh-lean/blob/72fead102f663d5c7a407964a13dc0cf9d977207/grh-compact/src/New/Data/WeilDataOfE.lean\#L144-L145}{{\leanfont Parameters.\allowbreak{}frobElem\_eq\_one}}}{WeilDataOfE.lean}{144--145}
\crow{splitne}{\href{https://github.com/lsernau/grh-lean/blob/72fead102f663d5c7a407964a13dc0cf9d977207/grh-compact/src/New/Data/WeilDataOfE.lean\#L151-L152}{{\leanfont Parameters.\allowbreak{}smul\_ne\_star}}}{WeilDataOfE.lean}{151--152}
\crow{splitfix}{\href{https://github.com/lsernau/grh-lean/blob/72fead102f663d5c7a407964a13dc0cf9d977207/grh-compact/src/New/Data/WeilDataOfE.lean\#L154-L155}{{\leanfont Parameters.\allowbreak{}smul\_eq\_self\_iff\_eq\_one}}}{WeilDataOfE.lean}{154--155}
\crow{splitq}{\href{https://github.com/lsernau/grh-lean/blob/72fead102f663d5c7a407964a13dc0cf9d977207/grh-compact/src/New/Data/WeilDataOfE.lean\#L158-L158}{{\leanfont Parameters.\allowbreak{}q\_eq\_of\_modEq\_one}}}{WeilDataOfE.lean}{158--158}
\crow{starmem}{\href{https://github.com/lsernau/grh-lean/blob/72fead102f663d5c7a407964a13dc0cf9d977207/grh-compact/src/New/Data/WeilDataOfE.lean\#L164-L166}{{\leanfont Parameters.\allowbreak{}star\_mem\_iff}}}{WeilDataOfE.lean}{164--166}
\crow{valnorm}{\href{https://github.com/lsernau/grh-lean/blob/72fead102f663d5c7a407964a13dc0cf9d977207/grh-compact/src/New/Data/WeilDataOfE.lean\#L170-L172}{{\leanfont Parameters.\allowbreak{}intValuation\_eq\_of\_mul\_star}}}{WeilDataOfE.lean}{170--172}
\crow{bij}{\href{https://github.com/lsernau/grh-lean/blob/72fead102f663d5c7a407964a13dc0cf9d977207/grh-compact/src/New/Data/WeilDataOfE.lean\#L178-L183}{{\leanfont Parameters.\allowbreak{}nonempty\_infinitePlace\_equiv\_finPlace}}}{WeilDataOfE.lean}{178--183}
\crow{valmap}{\href{https://github.com/lsernau/grh-lean/blob/72fead102f663d5c7a407964a13dc0cf9d977207/grh-compact/src/New/Data/WeilDataOfE.lean\#L192-L196}{{\leanfont Parameters.\allowbreak{}valuation\_map\_relDifferentGen}}}{WeilDataOfE.lean}{192--196}
\crow{alP}{\href{https://github.com/lsernau/grh-lean/blob/72fead102f663d5c7a407964a13dc0cf9d977207/grh-compact/src/New/Data/WeilDataOfE.lean\#L200-L202}{{\leanfont Parameters.\allowbreak{}α}}}{WeilDataOfE.lean}{200--202}
\crow{alinv}{\href{https://github.com/lsernau/grh-lean/blob/72fead102f663d5c7a407964a13dc0cf9d977207/grh-compact/src/New/Data/WeilDataOfE.lean\#L207-L207}{{\leanfont Parameters.\allowbreak{}inv\_α\_mem}}}{WeilDataOfE.lean}{207--207}
\crow{valpha}{\href{https://github.com/lsernau/grh-lean/blob/72fead102f663d5c7a407964a13dc0cf9d977207/grh-compact/src/New/Data/WeilDataOfE.lean\#L211-L217}{{\leanfont Parameters.\allowbreak{}v\_alpha\_mul\_relDifferentGen}}}{WeilDataOfE.lean}{211--217}
\crow{Nal}{\href{https://github.com/lsernau/grh-lean/blob/72fead102f663d5c7a407964a13dc0cf9d977207/grh-compact/src/New/Data/WeilDataOfE.lean\#L219-L220}{{\leanfont Parameters.\allowbreak{}N₁\_mul\_α\_mem}}}{WeilDataOfE.lean}{219--220}
\crow{weil}{\href{https://github.com/lsernau/grh-lean/blob/72fead102f663d5c7a407964a13dc0cf9d977207/grh-compact/src/New/Data/WeilDataOfE.lean\#L225-L243}{{\leanfont Parameters.\allowbreak{}instWeilDataE}}}{WeilDataOfE.lean}{225--243}}
% the notations of the defined functions: one slot for each explicit argument of the Lean declaration
\newcommand{\twd}[1]{\operatorname{tw}(#1)}
\newcommand{\Nn}[2]{N_1(#1,#2)}
\newcommand{\KP}[1]{K(#1)}
\newcommand{\Mom}[2]{M_\omega(#1,#2)}
\newcommand{\EP}[2]{E(#1,#2)}
\newcommand{\PlP}[2]{\operatorname{Pl}(#1,#2)}
\newcommand{\rtu}[2]{\varsigma(#1,#2)}
\newcommand{\OmP}[2]{\Omega(#1,#2)}
\newcommand{\galP}[2]{\operatorname{gal}_{#1,#2}}
\newcommand{\cycP}[2]{\operatorname{cyc}_{#1,#2}}
\newcommand{\pin}[3]{[#3]_{#1,#2}}
\newcommand{\pcl}[2]{\operatorname{cl}_{#1}^{#2}}
\newcommand{\prm}[2]{\operatorname{perm}_{#1}(#2)}
\newcommand{\rti}[2]{\varsigma(#1,#2)}
\newcommand{\czE}[2]{c_0(#1,#2)}
\newcommand{\sqd}[2]{\sqrt{-d}_{#1,#2}}
\newcommand{\dgen}[2]{\partial(#1,#2)}
\newcommand{\vpi}[2]{\varpi(#1,#2)}
\newcommand{\galv}[1]{\operatorname{gal}_{#1}}
\newcommand{\Frb}[2]{\operatorname{Fr}_{#1}(#2)}
\newcommand{\alP}[2]{\alpha(#1,#2)}
\newcommand{\UnitMembers}{239}
\newcommand{\UnitInNote}{96}
\newcommand{\UnitProofs}{39}
\newcommand{\UnitSteps}{145}
\newcommand{\UnitAllSteps}{172}


\linespread{1.08}                       % 11 pt on 14.7 pt
\setlength{\parindent}{0pt}
\setlength{\parskip}{5.5pt plus 1.5pt minus 0.5pt}
\setlength{\abovedisplayskip}{7.5pt plus 2pt minus 3pt}
\setlength{\belowdisplayskip}{7.5pt plus 2pt minus 3pt}
\setlength{\abovedisplayshortskip}{3pt plus 2pt}
\setlength{\belowdisplayshortskip}{6pt plus 2pt minus 2pt}
\emergencystretch=1.5em
\clubpenalty=10000 \widowpenalty=10000 \displaywidowpenalty=10000
\relpenalty=7000 \binoppenalty=9000
\raggedbottom
\pagestyle{fancy}\fancyhf{}\renewcommand{\headrulewidth}{0pt}
\fancyfoot[C]{\sffamily\fontsize{7.6}{9}\selectfont\color{inkfaint}\thepage\,/\,\pageref*{LastPage}}
\setlist{topsep=2pt,partopsep=0pt,parsep=0pt,itemsep=4pt,leftmargin=2.1em,labelsep=0.55em}

\newcommand{\Q}{\mathbb{Q}}
\newcommand{\Z}{\mathbb{Z}}
\newcommand{\F}{\mathbb{F}}
\newcommand{\len}{\operatorname{len}}
\newcommand{\cok}{\operatorname{cok}}
\newcommand{\End}{\operatorname{End}}
\newcommand{\Hom}{\operatorname{Hom}}
\newcommand{\Gal}{\operatorname{Gal}}
\newcommand{\diag}{\operatorname{diag}}
\newcommand{\kw}[1]{\emph{#1}}
\newcommand{\m}[1]{\mbox{$#1$}}        % a short formula that must not be broken across lines
\newcommand{\head}[1]{\par\addvspace{17pt plus 4pt minus 2pt}\Needspace{6\baselineskip}%
  \noindent{\fontsize{13.2}{16}\selectfont\bfseries\color{ink}#1}\par\nopagebreak\vspace{-1pt}}
\newcommand{\sect}[2]{\refstepcounter{section}\label{#2}\par\addvspace{19pt plus 4pt minus 2pt}\Needspace{7\baselineskip}%
  \noindent\llap{{\fontsize{13.2}{16}\selectfont\bfseries\color{inkfaint}\thesection}\hspace{9pt}}%
  {\fontsize{13.2}{16}\selectfont\bfseries\color{ink}#1}\par\nopagebreak\vspace{-1pt}}
\newcommand{\subhead}[1]{\par\addvspace{9pt plus 2pt}\Needspace{4\baselineskip}\noindent\textbf{#1}\enspace\ignorespaces}
% where a statement lives in the Lean archive
\newcommand{\leantag}{{\sffamily\bfseries\fontsize{6.3}{8}\selectfont\color{leancol}\addfontfeatures{LetterSpace=7}LEAN}}
\newcommand{\sep}{\unskip\nobreak\hspace{5pt}{\color{inkfaint}\textperiodcentered}\hspace{5pt plus 2pt}\ignorespaces}
\newcommand{\leanline}[1]{\par\nopagebreak\vspace{-2.5pt}{\fontsize{8.9}{12.4}\selectfont\color{inksoft}\raggedright
  \noindent\llap{\leantag\hspace{9pt}}#1\par}\vspace{1pt}}
% a displayed statement in a tinted box: #1 colour, #2 text
\newcommand{\statement}[2]{\par\addvspace{7pt}\begingroup\setlength{\fboxsep}{8.5pt}\setlength{\fboxrule}{0.6pt}%
  \noindent\fcolorbox{#1!30!white}{#1!3!white}{\parbox{\dimexpr\linewidth-2\fboxsep-2\fboxrule\relax}{#2}}%
  \endgroup\par\addvspace{7pt}}
\newcommand{\nt}[2]{\par\noindent\hangindent=1.1em\hangafter=1 #1\enspace{\color{inksoft}#2}\par\vspace{1pt}}
\newcommand{\legdot}[1]{\tikz[baseline=-0.62ex]\fill[#1] (0,0) circle (2.5pt);\hspace{3pt}}
\newcommand{\crow}[4]{\csname CS@#1\endcsname & #2 & \textcolor{inksoft}{#3} & \textcolor{inksoft}{#4}\\}
\newcommand{\leanquoteline}[2]{\par\noindent\hangindent=\dimexpr#1\fontdimen2\font+2.2em\relax\hangafter=1
  \hspace*{#1\fontdimen2\font}#2}
\newcommand{\quoted}[1]{\par\vspace{3pt}{\ttfamily\fontsize{7.5}{10.4}\selectfont\color{leancol}\raggedright
  \setlength{\parskip}{0pt}\leanquote{#1}\par}\vspace{4pt}}
\newcommand{\bx}[3]{{\sffamily\bfseries\fontsize{7.3}{9.2}\selectfont\color{#1!82!black}#2}\\[1.5pt]#3}
\newcommand{\cs}[2]{\expandafter\newcommand\csname CS@#1\endcsname{#2}}

% theorems
\newtheoremstyle{plainx}{7pt plus 2pt}{7pt plus 2pt}{\itshape}{}{\bfseries}{.}{0.5em}{}
\newtheoremstyle{uprightx}{7pt plus 2pt}{7pt plus 2pt}{}{}{\bfseries}{.}{0.5em}{}
\theoremstyle{plainx}
\newtheorem{theorem}{Theorem}[section]
\newtheorem{lemma}[theorem]{Lemma}
\newtheorem{proposition}[theorem]{Proposition}
\newtheorem{corollary}[theorem]{Corollary}
\theoremstyle{uprightx}
\newtheorem{definition}[theorem]{Definition}
\newtheorem{setting}[theorem]{Setting}          % hypotheses shared by several statements; each of them cites it by number
\newtheorem{remark}[theorem]{Remark}
\newtheorem{example}[theorem]{Example}
\renewcommand{\qedsymbol}{{\setlength{\fboxsep}{0pt}\setlength{\fboxrule}{0.5pt}\fbox{\rule{0pt}{0.46em}\hspace{0.46em}}}}
% what holds throughout the note: the ambient objects only (implicit and instance arguments), never a hypothesis
\newenvironment{conventions}{\par\addvspace{5pt}\noindent\textbf{Conventions.}\enspace\ignorespaces}{\par\addvspace{3pt}}
\makeatletter
\renewenvironment{proof}[1][Proof]{\par\pushQED{\qed}\normalfont\topsep0pt\partopsep0pt
  \trivlist\item[\hskip\labelsep\itshape #1.]\ignorespaces}{\popQED\endtrivlist\@endpefalse}
\makeatother

% where each declaration of the concordance is treated: \cs{key}{place}, the place naming a \label or a lettered item
\cs{par}{Definition~\ref{d:par}}
\cs{twd}{Definition~\ref{d:tw}}
\cs{Nn}{Definition~\ref{d:tw}}
\cs{KP}{Definition~\ref{d:tw}}
\cs{Mom}{Definition~\ref{d:tw}}
\cs{heeg}{Lemma~\ref{l:arith}}
\cs{ltd}{Lemma~\ref{l:arith}}
\cs{dvdnorm}{Lemma~\ref{l:arith}}
\cs{chione}{Lemma~\ref{l:arith}}
\cs{thirteen}{Lemma~\ref{l:arith}}
\cs{NdM}{Lemma~\ref{l:arith}}
\cs{coprime}{Lemma~\ref{l:arith}}
\cs{twopow}{Lemma~\ref{l:arith}}
\cs{powd}{Lemma~\ref{l:powd}}
\cs{EP}{Definition~\ref{d:E}}
\cs{PlP}{Definition~\ref{d:E}}
\cs{rtu}{Definition~\ref{d:E}}
\cs{OmP}{Definition~\ref{d:E}}
\cs{primrt}{Definition~\ref{d:E}}
\cs{irr}{Lemma~\ref{l:irr}}
\cs{tocyc}{Lemma~\ref{l:cyclo}}
\cs{embeq}{Lemma~\ref{l:emb}}
\cs{galP}{Definition~\ref{d:group}}
\cs{cycP}{Definition~\ref{d:group}}
\cs{pin}{Definition~\ref{d:group}}
\cs{pcl}{Definition~\ref{d:group}}
\cs{phialg}{Lemma~\ref{l:dict}}
\cs{phimul}{Lemma~\ref{l:dict}}
\cs{starrt}{Lemma~\ref{l:dict}}
\cs{galrt}{Theorem A}
\cs{pinpow}{Lemma~\ref{l:dict}}
\cs{clpin}{Lemma~\ref{l:dict}}
\cs{galpin}{Theorem A}
\cs{clneg}{Lemma~\ref{l:dict}}
\cs{expandpow}{Theorem A}
\cs{expandtw}{Theorem A}
\cs{cardPl}{Theorem A}
\cs{cardcl}{Theorem A}
\cs{prm}{Definition~\ref{d:perm}}
\cs{prmspec}{Lemma~\ref{l:perm}}
\cs{roots}{Lemma~\ref{l:roots}}
\cs{prodphi}{Lemma~\ref{l:roots}}
\cs{rti}{Definition~\ref{d:ints}}
\cs{czE}{Definition~\ref{d:ints}}
\cs{sqd}{Definition~\ref{d:ints}}
\cs{dgen}{Definition~\ref{d:ints}}
\cs{coert}{Definition~\ref{d:ints}}
\cs{primrti}{Definition~\ref{d:ints}}
\cs{coesq}{Definition~\ref{d:ints}}
\cs{starsq}{Lemma~\ref{l:ints}}
\cs{starrtpow}{Lemma~\ref{l:ints}}
\cs{sqdvd}{Lemma~\ref{l:ints}}
\cs{czmul}{Lemma~\ref{l:ints}}
\cs{starcz}{Lemma~\ref{l:ints}}
\cs{relspec}{Theorem C}
\cs{dgendvd}{Lemma~\ref{l:dif}}
\cs{dgenne}{Lemma~\ref{l:dif}}
\cs{difmul}{Theorem C}
\cs{condmem}{Lemma~\ref{l:dif}}
\cs{unitreal}{Lemma~\ref{l:real}}
\cs{vpi}{Definition~\ref{d:varpi}}
\cs{starvpi}{Lemma~\ref{l:varpi}}
\cs{difvpi}{Theorem C}
\cs{prodvpi}{Theorem C}
\cs{absvpi}{Lemma~\ref{l:varpi}}
\cs{galv}{Definition~\ref{d:frob}}
\cs{coegal}{Definition~\ref{d:frob}}
\cs{memsmul}{Definition~\ref{d:frob}}
\cs{Frb}{Definition~\ref{d:frob}}
\cs{copNz}{Lemma~\ref{l:frobext}}
\cs{frobext}{Lemma~\ref{l:frobext}}
\cs{frobpow}{Lemma~\ref{l:frobpow}}
\cs{decomp}{Theorem B}
\cs{qpow}{Theorem B}
\cs{notsq}{Lemma~\ref{l:unram}}
\cs{valp}{Lemma~\ref{l:unram}}
\cs{transit}{Lemma~\ref{l:trans}}
\cs{inertnd}{Lemma~\ref{l:inert}}
\cs{inertfix}{Lemma~\ref{l:inert}}
\cs{inertstar}{Lemma~\ref{l:inert}}
\cs{inertq}{Lemma~\ref{l:inert}}
\cs{splitchi}{Lemma~\ref{l:split}}
\cs{splitnd}{Lemma~\ref{l:split}}
\cs{splitfr}{Lemma~\ref{l:split}}
\cs{splitne}{Lemma~\ref{l:split}}
\cs{splitfix}{Lemma~\ref{l:split}}
\cs{splitq}{Lemma~\ref{l:split}}
\cs{starmem}{Lemma~\ref{l:norm}}
\cs{valnorm}{Lemma~\ref{l:norm}}
\cs{bij}{Proposition~\ref{p:bij}}
\cs{valmap}{Lemma~\ref{l:valmap}}
\cs{alP}{Definition~\ref{d:alpha}}
\cs{alinv}{Lemma~\ref{l:alpha}}
\cs{valpha}{Lemma~\ref{l:alpha}}
\cs{Nal}{Lemma~\ref{l:alpha}}
\cs{weil}{Proposition~\ref{p:weil}}

% ------------------------------------------------------------------ macros of this note that are not notations of
% defined functions (those are made by the build, in data.tex, from the spec)
\newcommand{\OO}{\mathfrak{o}}
\newcommand{\CC}{\mathbb{C}}
\newcommand{\RR}{\mathbb{R}}
\newcommand{\NN}{\mathbb{N}}
\newcommand{\Dif}{\mathfrak{D}}
\newcommand{\Cy}[1]{\Psi_{#1}}                 % the cyclotomic polynomial
\newcommand{\chiK}[1]{\chi_{\KP{#1}}}         % the quadratic character of K(P)
\newcommand{\emb}[1]{\iota_{#1}}              % the chosen embedding at an infinite place
\newcommand{\cj}[1]{\overline{#1}}
\newcommand{\ord}{\operatorname{ord}}
\newcommand{\Nm}{\operatorname{N}}
\newcommand{\Aut}{\operatorname{Aut}}
\newcommand{\GL}{\operatorname{GL}}
% components of the data D (fields of the record; not declarations of their own)
\newcommand{\Dm}[1]{m(#1)}
\newcommand{\Dr}[1]{r(#1)}
\newcommand{\Dmu}[1]{\mu(#1)}
\newcommand{\thick}[1]{\operatorname{thick}(#1)}
\newcommand{\Dg}[2]{g_{#1}(#2)}
\newcommand{\Dpig}[1]{\pi_g(#1)}
\newcommand{\ordch}[1]{\operatorname{ord}_{\mathrm{ch}}(#1)}
\newcommand{\Kch}[2]{K_{\mathrm{ch}}(#1,#2)}
\newcommand{\rdE}[2]{\operatorname{rd}_{\EP{#1}{#2}}}

\begin{document}
\color{ink}

{\raggedright
{\sffamily\bfseries\fontsize{7.2}{9}\selectfont\color{leancol}\addfontfeatures{LetterSpace=8}READBACK NOTE 5 \textperiodcentered{} CHAPTER III, FIRST PART\par}
\vspace{5pt}
{\fontsize{21.5}{25.5}\selectfont\bfseries The parameters and the CM field: places, Frobenius elements, the different\par}
\vspace{5pt}
{\fontsize{12.4}{16}\selectfont\itshape\color{inksoft} One unit of the Lean artifact \ArtifactName, read back into
mathematics\par}}
\vspace{8pt}
{\color{rulegrey}\hrule height 0.6pt}
\vspace{7pt}
\thispagestyle{empty}

\textbf{Summary.} The argument works over one family of number fields. The parameters $P=(N_0,d)$ are a multiple
$N_0$ of $18480$ and a prime $d\equiv N_0/2-1\bmod N_0$; they give the imaginary quadratic field
$\KP{P}=\Q(\sqrt{-d})$ and, for every exponent $\mu$, the CM field $\EP{P}{\mu}$ obtained by adjoining a root of
unity of order $\Nn{P}{\mu}=N_0\cdot3^{\lfloor\log_3\mu\rfloor+5}\cdot2^{\mu}$. This note sets up the arithmetic of
these fields: the infinite places of $\EP{P}{\mu}$ form a group, identified with its Galois group over $\KP{P}$ and
with the units modulo $\Nn{P}{\mu}$ (Theorem~A); the decomposition of a prime is read from an explicit Frobenius
element (Theorem~B); the ring of integers is generated by the root of unity and $(1+\sqrt{-d})/2$, and the different
is principal, generated by a purely imaginary integer the geometric mean of whose absolute values at the infinite
places is the root discriminant (Theorem~C). From these the note builds the data that the Weil representation over $\EP{P}{\mu}$ needs.

\statement{leancol}{\fontsize{9.6}{13.2}\selectfont\textbf{What this note is.} The artifact \ArtifactName{} is a
body of Lean code whose final theorems state the Riemann hypothesis and its analogue for Dirichlet $L$-functions.
That is a claim: its authors report that Lean's kernel accepts the code; a replay of that check by the owner of
this write-up was started, and the handout's record of checks gives its state; nothing was compiled for this note;
no mathematician has read the proof through. This note treats one part of one unit (the folder of the data of the proof): \UnitInNote{} of its
\UnitMembers{} declarations, those of the two files on the parameters, the CM field and its Weil data; the note
\emph{Admissible data} treats the rest. Every statement below was compared, argument by argument, with the Lean
text, and every proof is given in full, except that results of the background library are quoted as \emph{Inputs}
with their place in the artifact. The Lean text read is \PinText; Appendix~A is the concordance of \ConcDecls{}
declarations.}

\sect{The statements}{s:intro}
The notation is that of Section~\ref{s:setting}: $P=(N_0,d)$ are parameters, $\mu$ is a natural number,
$\Nn{P}{\mu}=N_0\,\twd{\mu}$ is the order of the root of unity $\rtu{P}{\mu}$ of the field $\EP{P}{\mu}$ over
$\KP{P}=\Q(\sqrt{-d})$, and $\PlP{P}{\mu}$ is the set of infinite places of $\EP{P}{\mu}$, a group by
Definition~\ref{d:group}.

\statement{arith}{\phantomsection\label{t:A}\textbf{Theorem A (places, automorphisms, classes).}\enspace Let $P=(N_0,d)$
be parameters and $\mu$ a natural number.
\begin{enumerate}[label=(\alph*)]
\item For every infinite place $\tau\in\PlP{P}{\mu}$, the automorphism $\galP{P}{\mu}(\tau)$ of $\EP{P}{\mu}$ over $\KP{P}$
sends $\rtu{P}{\mu}$ to $\rtu{P}{\mu}^{\,n}$, where $n\in\{0,\dots,\Nn{P}{\mu}-1\}$ represents the class
$\cycP{P}{\mu}(\tau)\in(\Z/\Nn{P}{\mu})^\times$.
\item For every integer $\nu$ prime to $N_0$, the automorphism $\galP{P}{\mu}(\pin{P}{\mu}{\nu})$ sends
$\rtu{P}{\mu}$ to $\rtu{P}{\mu}^{\,\nu}$.
\item $\EP{P}{\mu}$ has exactly $\twd{\mu}\,\varphi(N_0)$ infinite places, $\varphi$ being Euler's function.
\item Exactly $\twd{\mu}$ places $\tau\in\PlP{P}{\mu}$ have class $1$ modulo $N_0$, that is
$\pcl{P}{N_0\mid\Nn{P}{\mu}}(\tau)=1$.
\end{enumerate}}

\statement{arith}{\phantomsection\label{t:B}\textbf{Theorem B (decomposition of primes).}\enspace Let $P=(N_0,d)$ be
parameters, $\mu$ a natural number, $p$ a prime number with $p\nmid\Nn{P}{\mu}\,d$, and $w$ a finite place of
$\EP{P}{\mu}$ with $p\in w$. Let $\tau\in\PlP{P}{\mu}$, and let $(\mathfrak{c},e)$ be one of the two pairs
$(\mathrm{id},1)$ and $(\text{complex conjugation},-1)$, where $\mathfrak{c}$ is an automorphism of the ring of
integers of $\EP{P}{\mu}$ and $e$ an integer.
\begin{enumerate}[label=(\alph*)]
\item $\tau\cdot w=\mathfrak{c}^{-1}(w)$ if and only if there is a natural number $n$ with
$(\tau,e)=\Frb{\mu}{p}^{\,n}$ in $\PlP{P}{\mu}\times\Z$. Here $\mathfrak{c}^{-1}(w)$ is $w$ for the first pair and the
conjugate place $\cj{w}$ for the second.
\item The norm of $w$ is $p^{\,o}$, where $o$ is the order of the Frobenius element $\Frb{\mu}{p}$ in
$\PlP{P}{\mu}\times\Z$.
\end{enumerate}}

\statement{arith}{\phantomsection\label{t:C}\textbf{Theorem C (integers and the different).}\enspace Let $P=(N_0,d)$ be
parameters and $\mu$ a natural number; write $\OO$ for the ring of integers of $\EP{P}{\mu}$ and $\OO_K$ for that of
$\KP{P}$.
\begin{enumerate}[label=(\alph*)]
\item Let $z\in\OO$ be a primitive root of unity of order $\Nn{P}{\mu}$ and $y=\Cy{\Nn{P}{\mu}}'(z)$ the value at $z$
of the derivative of the cyclotomic polynomial of that order. Then $y$ divides $\Nn{P}{\mu}$ in $\OO$, the ideal
$(y)$ is the different of $\OO$ over $\OO_K$, and $\OO=\OO_K[z]$.
\item The different of $\OO$ over $\Z$ is the principal ideal generated by $\dgen{P}{\mu}\cdot\sqd{P}{\mu}$, and
also by the purely imaginary integer $\vpi{P}{\mu}$ of Definition~\ref{d:varpi}.
\item $\prod_{\tau\in\PlP{P}{\mu}}|\emb{\tau}(\vpi{P}{\mu})|=\rdE{P}{\mu}^{\,\#\PlP{P}{\mu}}$, where $\rdE{P}{\mu}$ is the root discriminant of $\EP{P}{\mu}$.
\end{enumerate}}

\textbf{Why these parameters.} The explanations of this paragraph are motivation, read from the comments of the
formal text and from the places where later units use the results; they are not statements of this unit. The
congruence on $d$ does two things. First, $d\equiv-1$ modulo $N_0/2$ and $8\mid N_0/2$ make every prime $p$ that
divides $\Nn{P}{\mu}$ a divisor of $(d+1)/4$, the norm of $(1+\sqrt{-d})/2$. So all these primes split in $\KP{P}$,
into the prime above $p$ that contains $(1+\sqrt{-d})/2$ and its conjugate, and the different of $\EP{P}{\mu}$ over
$\KP{P}$, which lives above these primes, can be divided between an idele and its conjugate
(Lemma~\ref{l:alpha}); this is the self-duality that the Weil representation asks for. Second,
$d\not\equiv-1$ modulo $N_0$ itself makes the powers of $d$ modulo $N_0$ avoid $-1$ (Lemma~\ref{l:powd}). For a
natural number $n$ and a place $\tau$, the place $\pin{P}{\mu}{-d^{\,n}}\cdot \tau$ then never has the class of $\tau$
modulo $N_0$: the ramified prime $d$ never relates two places of the same class, and later units use this. (The
prime $d$ does occur in the prime sums of the main inequality: $d^2\equiv1\bmod N_0$ by the proof of
Lemma~\ref{l:powd}.) The divisibility by $18480=2^4\cdot3\cdot5\cdot7\cdot11$ makes every prime below $13$ split in
$\KP{P}$ (Lemma~\ref{l:arith}(c)); with the factor $3^5$ of $\twd{\mu}$ it puts the roots of unity of order
$5040=2^4\cdot3^2\cdot5\cdot7$ into $\EP{P}{\mu}$. The order of the root of unity grows with $\mu$ only by powers of $2$ and $3$, so that the count of
places is exact (Theorem~A(c), (d)) and the group of places has a large cyclic $2$-part, where the note
\emph{Admissible data} builds the thick places.

\sect{The setting}{s:setting}
\begin{conventions}\phantomsection\label{c:conventions}
$\NN=\{0,1,2,\dots\}$. For a number field $F$, a \emph{finite place} is a nonzero prime ideal of its ring of
integers, and an \emph{infinite place} is an archimedean absolute value; the place of an embedding
$\psi\colon F\to\CC$ is $x\mapsto|\psi(x)|$. A CM field carries its complex conjugation $x\mapsto\cj{x}$, which
every embedding into $\CC$ turns into complex conjugation; it acts on the finite places, $\cj{w}=\{\cj{x}:x\in w\}$.
$\Cy{M}$ is the cyclotomic polynomial of order $M$, the monic integral polynomial whose complex roots are the
primitive roots of unity of order $M$. Three letters are local to this note: $\psi$, $\psi'$ are embeddings of
fields and homomorphisms of rings, and $\Psi$ occurs only in $\Cy{M}$; the columns $\psi$ of the overview of the
handout and the additive character $\Psi$ of Chapter~IV do not occur here. In a statement about a commutative ring $R$ of characteristic $p$, or an
integral domain $R$, these properties of $R$ are part of the statement. Divisibility, floors and subtraction of
natural numbers are those of $\NN$, as in the formal text. The letter $\nu$ is an integer, never a character of the ideles
(Convention~0.6), and from Section~\ref{s:field} on a bare $x$ is an element of a field or of a ring: this note has no scales
of data (Convention~0.3).
\end{conventions}

\begin{definition}[Parameters]\label{d:par}
\emph{Parameters} are a pair $P=(N_0,d)$ of natural numbers such that $18480$ divides $N_0$, $d$ is a prime number,
$N_0$ divides $2(d+1)$, and $N_0$ does not divide $d+1$. Since $N_0$ is even, the last two conditions say together
that $d\equiv N_0/2-1\bmod N_0$.
\end{definition}
\leanline{\lean{par}}

\begin{definition}[The order of the root of unity]\label{d:tw}
For a natural number $\mu$ put $\twd{\mu}=3^{\lfloor\log_3\mu\rfloor+5}\cdot2^{\mu}$, where
$\lfloor\log_3 0\rfloor=0$. For parameters $P=(N_0,d)$ and a natural number $\mu$ put
\[\Nn{P}{\mu}=N_0\cdot\twd{\mu},\qquad\Mom{P}{\mu}=(N_0\,d)^{\mu+6},\]
and let $\KP{P}=\Q(\mathrm{i}\sqrt d)\subset\CC$ be the imaginary quadratic field $\Q(\sqrt{-d})$, as a subfield of the
complex numbers. Its quadratic character $\chiK{P}$ is, on the integers, the Legendre symbol modulo $d$ (so
$\chiK{P}(d)=0$); in the formal text it is a function of $d$.
\end{definition}
\leanline{\lean{twd}\sep\lean{Nn}\sep\lean{KP}\sep\lean{Mom}}

The power of $3$ in $\twd{\mu}$ lies between $81\mu$ and $243\mu$ for $\mu\ge1$. The number $\Mom{P}{\mu}$ will be the
order of certain characters of residue fields in later units; here only Lemma~\ref{l:arith}(d) concerns it.

\begin{definition}[The CM field, its places, its root of unity]\label{d:E}
Let $P=(N_0,d)$ be parameters and $\mu$ a natural number. $\EP{P}{\mu}$ is the splitting field over $\KP{P}$ of
the cyclotomic polynomial $\Cy{\Nn{P}{\mu}}$ (in it $X^{\Nn{P}{\mu}}-1$ splits as well), and $\rtu{P}{\mu}\in\EP{P}{\mu}$ is a chosen primitive root of unity of order
$\Nn{P}{\mu}$; so $\EP{P}{\mu}=\KP{P}(\rtu{P}{\mu})$. $\PlP{P}{\mu}$ is the set of infinite places of $\EP{P}{\mu}$.
Finally
\[\OmP{P}{\mu}=\exp\bigl(\mathrm{cs}(d)\bigr)\cdot\frac{2\pi}{\rdE{P}{\mu}},\]
where $\rdE{P}{\mu}$ is the root discriminant of $\EP{P}{\mu}$ (the absolute value of its discriminant, raised to
the power one over its degree) and $\mathrm{cs}(d)$ is Euler's constant minus the real part of the logarithmic derivative
at $1$ of the entire Dirichlet series of $\chiK{P}$ (its $L$-function made entire) (a definition of the library, unit~106).
\end{definition}
\leanline{\lean{EP}\sep\lean{PlP}\sep\lean{rtu}\sep\lean{primrt}\sep\lean{OmP}}

The number $\OmP{P}{\mu}$ is a period: later units measure theta functions against it. This unit only defines it.

\sect{The parameters}{s:par}
\begin{lemma}[Arithmetic of the parameters]\label{l:arith}
Let $P=(N_0,d)$ be parameters and $\mu$ a natural number.
\begin{enumerate}[label=(\alph*)]
\item $d\equiv7\bmod8$, and $64<d$.
\item If $p$ is a prime number dividing $\Nn{P}{\mu}$, then $p$ divides $(d+1)/4$, and $\chiK{P}(p)=1$.
\item If $p$ is a prime number with $\chiK{P}(p)\ne1$, then $13\le p$.
\item $\Nn{P}{\mu}$ and $d$ are coprime; $2^{\mu+1}$ divides $\Nn{P}{\mu}$; and $\Nn{P}{\mu}\,d$ divides
$\Mom{P}{\mu}$.
\end{enumerate}
\end{lemma}
\begin{proof}
(a) $16$ divides $N_0$ and $N_0$ divides $2(d+1)$, so $8\mid d+1$. Also $N_0\ne0$, because $0$ divides only $0$, so
$18480\le N_0\le2(d+1)$.

(b) The primes $2$ and $3$ divide $N_0$, so $p$ divides $N_0$, hence $p$ divides $2(d+1)=8\cdot\frac{d+1}4$ (the
quotient is exact by~(a)). Thus $p$ divides
$8$ or $(d+1)/4$; in the first case $p=2$, and $2$ divides $(d+1)/4$ by~(a). So $4p$ divides $d+1$. Now
$\chiK{P}(p)$ is the Jacobi symbol of $p$ over the odd number $d$, and a Jacobi symbol with upper entry $p$ depends
only on the lower entry modulo $4p$. As $d\equiv4p-1\bmod4p$, it equals the symbol of $p$ over $4p-1$, which is the
symbol of $4p$ over $4p-1$ because $4$ is a square, that is the symbol of $1$ over $4p-1$, which is $1$.

(c) The primes below $13$ divide $18480$, hence $\Nn{P}{0}$, and (b) applies with the exponent $0$.

(d) If the prime $d$ divided $\Nn{P}{\mu}$, then $\chiK{P}(d)=1$ by~(b), but $\chiK{P}(d)=0$. Next $2\mid N_0$ and
$2^\mu\mid\twd{\mu}$. Finally $\lfloor\log_3\mu\rfloor\le\mu$, so $\twd{\mu}$ divides $6^{\mu+5}$, which divides
$(N_0d)^{\mu+5}$ because $6\mid N_0$; multiplying by $N_0\,d$ gives the claim.
\end{proof}
\leanline{\lean{heeg}\sep\lean{ltd}\sep\lean{dvdnorm}\sep\lean{chione}\sep\lean{thirteen}\sep\lean{coprime}\sep
\lean{twopow}\sep\lean{NdM}}

\begin{lemma}[Powers of $d$ modulo $N_0$]\label{l:powd}
Let $P=(N_0,d)$ be parameters and $n$ a natural number. Then $d^{\,n}\not\equiv-1\bmod N_0$.
\end{lemma}
\begin{proof}
In $\Z/N_0$ one has $2(d+1)=0$ and $d+1\ne0$. As $d$ is odd, $d-1=2i$ with an integer $i$, and
$d^2-1=i\cdot2(d+1)=0$. So every power of $d$ is $1$ or $d$. Here $d\ne-1$ by assumption, and $1\ne-1$, since
$2=0$ would give $d+1=2(i+1)=0$.
\end{proof}
\leanline{\lean{powd}}

\sect{The field and its places}{s:field}
\begin{lemma}[The cyclotomic polynomial stays irreducible]\label{l:irr}
Let $P=(N_0,d)$ be parameters and $\mu$ a natural number. Then $\Cy{\Nn{P}{\mu}}$ is irreducible over $\KP{P}$.
\end{lemma}
\begin{proof}
Let $F$ be the cyclotomic field of order $d$. As $d$ is a prime with $d\equiv3\bmod4$ (Lemma~\ref{l:arith}(a)),
$\KP{P}$ embeds into $F$ by Input~(I0) of Section~\ref{s:dif}. It suffices that the polynomial is irreducible over
$F$. Let $x$ be a root of it in an extension of $F$. The minimal polynomial of $x$ over $F$ divides
$\Cy{\Nn{P}{\mu}}$, and its degree is $[F(x):F]$. The field $F(x)$ contains roots of unity of the coprime orders
$\Nn{P}{\mu}$ and $d$ (Lemma~\ref{l:arith}(d)), hence of order $\Nn{P}{\mu}d$, so its degree over $\Q$ is at least
$\varphi(d)\,\varphi(\Nn{P}{\mu})$ by the irreducibility of cyclotomic polynomials over $\Q$. As $[F:\Q]=\varphi(d)$,
we get $[F(x):F]\ge\varphi(\Nn{P}{\mu})$, the degree of $\Cy{\Nn{P}{\mu}}$. So $\Cy{\Nn{P}{\mu}}$ is the minimal
polynomial of $x$.
\end{proof}
\leanline{\lean{irr}}

\begin{lemma}[The CM field lies in a cyclotomic field]\label{l:cyclo}
Let $P=(N_0,d)$ be parameters and $\mu$ a natural number. There is an embedding of fields of $\EP{P}{\mu}$ into the
cyclotomic field of order $\Nn{P}{\mu}\,d$. Consequently $\EP{P}{\mu}$ is abelian over $\Q$ and over $\KP{P}$, it is
totally complex, and it is a CM field.
\end{lemma}
\begin{proof}
The cyclotomic field of order $\Nn{P}{\mu}d$ contains a primitive root of unity of order $d$, hence a copy of
$\KP{P}$ by (I0), and over it $\Cy{\Nn{P}{\mu}}$ splits; a splitting field embeds into every
extension in which the polynomial splits. A subfield of an abelian extension of $\Q$ is abelian over $\Q$, and
$\EP{P}{\mu}$ is totally complex because it contains the imaginary quadratic field $\KP{P}$. A totally complex
abelian number field is a CM field: because the Galois group is abelian, complex conjugation is the same
automorphism for every embedding into $\CC$, and its fixed field is totally real.
\end{proof}
\leanline{\lean{tocyc}; the consequences are instances without names, stated after it in the formal text.}

\begin{lemma}[An embedding over the quadratic field is determined by its place]\label{l:emb}
Let $P=(N_0,d)$ be parameters, $\mu$ a natural number, and $\psi$, $\psi'$ two embeddings of $\EP{P}{\mu}$ into
$\CC$ that are the identity on $\KP{P}$. If $\psi$ and $\psi'$ have the same infinite place, then $\psi=\psi'$.
\end{lemma}
\begin{proof}
Two embeddings with the same place are equal or complex conjugate. If $\psi'$ were the conjugate of $\psi$, then
every $x\in\KP{P}$ would be equal to its complex conjugate, which fails for $\mathrm{i}\sqrt d$.
\end{proof}
\leanline{\lean{embeq}}

\begin{definition}[The group of places and its three descriptions]\label{d:group}
Let $P=(N_0,d)$ be parameters and $\mu$ a natural number. Fix an embedding $\emb{1}\colon\EP{P}{\mu}\to\CC$ that is
the identity on $\KP{P}$ (the formal text takes the one that the algebraic closedness of $\CC$ provides). The map
\[\Gal(\EP{P}{\mu}/\KP{P})\longrightarrow\PlP{P}{\mu},\qquad \gamma\longmapsto\text{the place of }\emb{1}\circ\gamma,\]
is a bijection: it is injective by Lemma~\ref{l:emb}, and surjective because the Galois group permutes
transitively the places of $\EP{P}{\mu}$ above the one infinite place of $\KP{P}$. The set $\PlP{P}{\mu}$ is given
the structure of a commutative group that makes this bijection an isomorphism; its unit element is the place
of~$\emb{1}$.
\begin{enumerate}[label=(\roman*)]
\item $\galP{P}{\mu}\colon\PlP{P}{\mu}\to\Gal(\EP{P}{\mu}/\KP{P})$ is the inverse isomorphism.
\item $\cycP{P}{\mu}\colon\PlP{P}{\mu}\to(\Z/\Nn{P}{\mu})^\times$ is $\galP{P}{\mu}$ followed by the cyclotomic
character, which is an isomorphism by Lemma~\ref{l:irr}.
\item For an integer $\nu$, the place $\pin{P}{\mu}{\nu}\in\PlP{P}{\mu}$ is the one with
$\cycP{P}{\mu}(\pin{P}{\mu}{\nu})=\nu\bmod\Nn{P}{\mu}$ if $\nu$ is prime to $\Nn{P}{\mu}$, and the unit element
otherwise.
\item For a natural number $m'$ with $m'\mid\Nn{P}{\mu}$, the homomorphism
\[\pcl{P}{m'\mid\Nn{P}{\mu}}\colon\PlP{P}{\mu}\longrightarrow(\Z/m')^\times\]
is $\cycP{P}{\mu}$ followed by reduction modulo $m'$: the \emph{class modulo $m'$} of a place. (In the formal text its explicit arguments are $P$ and the proof of the
divisibility; the divisibility is written in the place of that proof.)
\item For $\tau\in\PlP{P}{\mu}$ the \emph{chosen embedding at $\tau$} is $\emb{\tau}=\emb{1}\circ\galP{P}{\mu}(\tau)$; its place
is $\tau$. This is the CM type of $\EP{P}{\mu}$ that all later units use.
\end{enumerate}
\end{definition}
\leanline{\lean{galP}\sep\lean{cycP}\sep\lean{pin}\sep\lean{pcl}; the group structure and the CM type are instances
without names.}

\begin{lemma}[The dictionary]\label{l:dict}
Let $P=(N_0,d)$ be parameters and $\mu$ a natural number; let $\tau,\tau'\in\PlP{P}{\mu}$, let $\nu$ be an integer and
$i$ a natural number.
\begin{enumerate}[label=(\alph*)]
\item $\emb{\tau}(x)=x$ for every $x\in\KP{P}$, and $\emb{\tau\,\tau'}(x)=\emb{\tau}\bigl(\galP{P}{\mu}(\tau')(x)\bigr)$ for every
$x\in\EP{P}{\mu}$.
\item $\cj{\rtu{P}{\mu}}=\rtu{P}{\mu}^{-1}$.
\item $\pin{P}{\mu}{\nu}^{\,i}=\pin{P}{\mu}{\nu^i}$.
\item If $\nu$ is prime to $N_0$, then for every natural number $m'$ with $m'\mid\Nn{P}{\mu}$ the class
$\pcl{P}{m'\mid\Nn{P}{\mu}}(\pin{P}{\mu}{\nu})$ is the class of $\nu$ modulo $m'$.
\item $\pcl{P}{N_0\mid\Nn{P}{\mu}}(\pin{P}{\mu}{\nu})=-1$ if and only if $\nu\equiv-1\bmod N_0$.
\end{enumerate}
\end{lemma}
\begin{proof}
(a) holds by the definition of $\emb{\tau}$, since $\galP{P}{\mu}$ is a homomorphism into a group of automorphisms over
$\KP{P}$. (b) The injective map $\emb{1}$ commutes with conjugation and sends $\rtu{P}{\mu}$ to a complex number of
absolute value $1$, whose conjugate is its inverse. (c) If $\nu$ is prime to $\Nn{P}{\mu}$, both sides have the
class $\nu^i$. If not, then $\nu^i$ is not prime to $\Nn{P}{\mu}$ for $i\ge1$ and both sides are the unit element,
and for $i=0$ both sides are the unit element as well. (d) The numbers $N_0$ and $\Nn{P}{\mu}$ have the same prime
divisors, so $\nu$ is prime to $\Nn{P}{\mu}$ and the class of $\pin{P}{\mu}{\nu}$ is $\nu$. (e) If the class is
$-1$, then $\nu$ is prime to $\Nn{P}{\mu}$, because the unit element has class $1\ne-1$ (Lemma~\ref{l:powd} with
$n=0$); then (d) applies. Conversely $\nu\equiv-1\bmod N_0$ is prime to $N_0$, and (d) applies.
\end{proof}
\leanline{\lean{phialg}\sep\lean{phimul}\sep\lean{starrt}\sep\lean{pinpow}\sep\lean{clpin}\sep\lean{clneg}}

\begin{proof}[Proof of Theorem A]
(a) is the definition of the cyclotomic character, and (b) follows from (a) and Lemma~\ref{l:dict}(d): the exponent
$n$ represents $\nu$ modulo $\Nn{P}{\mu}$, and $\rtu{P}{\mu}$ has that order. (c) By
Definition~\ref{d:group} the number of places is the order of $(\Z/\Nn{P}{\mu})^\times$, which is the degree of
$\Cy{\Nn{P}{\mu}}$. For a prime $p$ dividing $M$ one has $\Cy{M}(X^{p^{n_0}})=\Cy{Mp^{n_0}}(X)$ (induction on $n_0$); as $2$ and
$3$ divide $N_0$, this gives $\Cy{\Nn{P}{\mu}}(X)=\Cy{N_0}(X^{\twd{\mu}})$, of degree $\twd{\mu}\,\varphi(N_0)$.
(d) The places of class $1$ modulo $N_0$ form the kernel of the reduction from $(\Z/\Nn{P}{\mu})^\times$ onto
$(\Z/N_0)^\times$, whose order is the quotient of the two orders.
\end{proof}
\leanline{\lean{galrt}\sep\lean{galpin}\sep\lean{cardPl}\sep\lean{cardcl}\sep\lean{expandpow}\sep\lean{expandtw}}

\begin{definition}[The permutation of the places by an automorphism of $\CC$]\label{d:perm}
Let $P=(N_0,d)$ be parameters, $\mu$ a natural number and $\sigma$ a field automorphism of $\CC$ that is the
identity on $\KP{P}$. Then $\prm{\mu}{\sigma}$ is the permutation of $\PlP{P}{\mu}$ that multiplies every place, on
the left, by the place of the embedding $\sigma\circ\emb{1}$. (Its explicit arguments in the formal text are $\mu$ and
$\sigma$; the parameters are read from $\sigma$.)
\end{definition}
\leanline{\lean{prm}}

\begin{lemma}\label{l:perm}
Let $P=(N_0,d)$ be parameters, $\mu$ a natural number, $\sigma$ a field automorphism of $\CC$ that is the identity on
$\KP{P}$, $\tau\in\PlP{P}{\mu}$ and $x\in\EP{P}{\mu}$. Then
$\sigma\bigl(\emb{\tau}(x)\bigr)=\emb{\tau^\ast}(x)$ with $\tau^\ast=\prm{\mu}{\sigma}(\tau)$.
\end{lemma}
\begin{proof}
Let $\tau'$ be the place of $\sigma\circ\emb{1}$. The embeddings $\sigma\circ\emb{1}$ and $\emb{\tau'}$ are the identity on
$\KP{P}$ and have the same place, so they are equal (Lemma~\ref{l:emb}). Hence
$\sigma(\emb{\tau}(x))=\emb{\tau'}(\galP{P}{\mu}(\tau)(x))=\emb{\tau'\,\tau}(x)$ by Lemma~\ref{l:dict}(a).
\end{proof}
\leanline{\lean{prmspec}}

\begin{lemma}[Roots of unity, and the norm to the quadratic field]\label{l:roots}
Let $P=(N_0,d)$ be parameters, $\mu$ a natural number and $y\in\EP{P}{\mu}$.
\begin{enumerate}[label=(\alph*)]
\item If $y$ is a root of unity, then there is a natural number $n<\Nn{P}{\mu}$ with $\rtu{P}{\mu}^{\,n}=y$.
\item $\prod_{\tau\in\PlP{P}{\mu}}\emb{\tau}(y)$ is the norm of $y$ from $\EP{P}{\mu}$ to $\KP{P}\subset\CC$.
\end{enumerate}
\end{lemma}
\begin{proof}
(a) It suffices that $y^{\Nn{P}{\mu}}=1$. By Lemma~\ref{l:cyclo} the roots of unity of $\EP{P}{\mu}$ embed into those
of the cyclotomic field of the even order $\Nn{P}{\mu}d$, a group of that order; so $z=y^{\Nn{P}{\mu}}$ satisfies
$z^d=1$. If $z\ne1$, then $z$ has order $d$, and as in the proof of Lemma~\ref{l:irr} the degree of $\EP{P}{\mu}$
over $\Q$ is at least $\varphi(d)\,\varphi(\Nn{P}{\mu})$. That degree is twice the number of infinite places,
$2\,\varphi(\Nn{P}{\mu})$, so $d-1\le2$, against $d>64$. (b) The norm is the product of the conjugates
$\gamma(y)$ over the Galois group; apply $\emb{1}$, which is the identity on $\KP{P}$.
\end{proof}
\leanline{\lean{roots}\sep\lean{prodphi}}

\sect{The different}{s:dif}
\begin{definition}[Four integers]\label{d:ints}
Let $P=(N_0,d)$ be parameters and $\mu$ a natural number, and let $\OO$ be the ring of integers of $\EP{P}{\mu}$.
The root of unity $\rti{P}{\mu}$ lies in $\OO$: it is the root of unity $\rtu{P}{\mu}$ of Definition~\ref{d:E},
written $\rti{P}{\mu}$ when it is regarded as an element of $\OO$ (the formal text has two names for it). $\czE{P}{\mu}\in\OO$ is the integer $(1+\sqrt{-d})/2$ of $\KP{P}$, and
$\sqd{P}{\mu}=2\,\czE{P}{\mu}-1$ is the square root of $-d$ that lies in $\KP{P}\subset\CC$ on the positive
imaginary axis. Finally $\dgen{P}{\mu}=\Cy{\Nn{P}{\mu}}'(\rti{P}{\mu})\in\OO$.
\end{definition}
\leanline{\lean{rti}\sep\lean{czE}\sep\lean{sqd}\sep\lean{dgen}\sep\lean{coert}\sep\lean{primrti}\sep\lean{coesq}}

\begin{lemma}\label{l:ints}
Let $P=(N_0,d)$ be parameters and $\mu$, $n$ natural numbers. Then, in the ring of integers of $\EP{P}{\mu}$:
$\cj{\sqd{P}{\mu}}=-\sqd{P}{\mu}$, and $\sqd{P}{\mu}$ divides $d$;
$\czE{P}{\mu}\,(1-\czE{P}{\mu})=(d+1)/4$ and $\cj{\czE{P}{\mu}}=1-\czE{P}{\mu}$; and
$\cj{\rti{P}{\mu}}^{\,n}\,\rti{P}{\mu}^{\,n}=1$.
\end{lemma}
\begin{proof}
The square of $\sqd{P}{\mu}$ is $-d$, its conjugate is the other root $-\sqd{P}{\mu}$, and
$\czE{P}{\mu}=(1+\sqd{P}{\mu})/2$. The last identity is Lemma~\ref{l:dict}(b).
\end{proof}
\leanline{\lean{starsq}\sep\lean{sqdvd}\sep\lean{czmul}\sep\lean{starcz}\sep\lean{starrtpow}}

\statement{ana}{\textbf{Inputs.}\enspace In (I0) to (I1$'$), $d$ is any prime number with the congruence stated.
(I0), unit~106: for a prime $d\equiv3\bmod4$, the field $\Q(\sqrt{-d})$ embeds into every field of characteristic
$0$ that contains a primitive root of unity of order $d$ (a quadratic Gauss sum).
(I1), unit~106: for a prime $d\equiv7\bmod8$ the different of
$\Q(\sqrt{-d})$ over $\Q$ is generated by $\sqrt{-d}$, and the ring of integers is $\Z[(1+\sqrt{-d})/2]$.
(I1$'$), unit~106: let $d\equiv7\bmod8$ be prime, $p$ an odd prime number, $\psi'$ a homomorphism from the ring of
integers of $\Q(\sqrt{-d})$ to a commutative ring $R$ of characteristic $p$, $n$ a natural number and $x\in R$.
Then $x=\psi'\bigl((1+\sqrt{-d})/2\bigr)^{p^n}$ if and only if $2x-1=\chi(p)^n\,\psi'(\sqrt{-d})$, where $\chi(p)$ is
the Legendre symbol of $p$ modulo $d$ (Euler's criterion with the law of quadratic reciprocity).
(I2), unit~10: let $F\subset F'$ be number fields with rings of integers $\OO_F\subset\OO_{F'}$, let $z\in\OO_{F'}$
generate $F'$ over $F$, and let a polynomial over $\OO_F$ vanish at $z$ and have derivative $y$ at $z$; if the ideal
$(y)$ divides the different of $\OO_{F'}$ over $\OO_F$, then that different is $(y)$ and $\OO_{F'}=\OO_F[z]$.
(I3), unit~10: in the ring of integers of a number field, an ideal whose square contains a prime number divides the
different over~$\Z$; the ideal need not be prime (Lemma~\ref{l:unram} applies (I3) to a finite place).

\smallskip References, as the catalogue of the write-up and the dossier of the unit record them; none was opened
for this note. For (I3): K.~Conrad~\cite{conrad}, Theorem~4.13 and its proof, of which (I3) is a consequence (the
lower bound for the exponent of the different at a ramified prime; the reading agent's record), and
Neukirch~\cite{neukirch}, Ch.~III, \S2, place not confirmed. For (I1): Neukirch~\cite{neukirch}, Ch.~I, place not
confirmed. For (I2): Serre~\cite{serre}, Ch.~III, place not confirmed. For (I0) and (I1$'$) the catalogue names no
work. From Mathlib we use that the ring of integers of $\Q(z)$, for a root of unity $z$, is $\Z[z]$, and that the
different is multiplicative in a tower of number fields.}

\begin{proof}[Proof of Theorem C(a)]
Write $X^{\Nn{P}{\mu}}-1=\Cy{\Nn{P}{\mu}}\cdot Q$ with an integral polynomial $Q$. Differentiating and evaluating
at $z$ gives $y\,Q(z)=\Nn{P}{\mu}\,z^{\Nn{P}{\mu}-1}$, so $y\,Q(z)\,z=\Nn{P}{\mu}$: thus $y$ divides
$\Nn{P}{\mu}$. The subfield $\Q(z)$ is the cyclotomic field of order $\Nn{P}{\mu}$; its ring of integers is $\Z[z]$,
so its different over $\Q$ is generated by $y$. By (I1) the different of $\KP{P}$ over $\Q$ is generated by
$\sqd{P}{\mu}$, which divides $d$ (Lemma~\ref{l:ints}); as $\Nn{P}{\mu}$ and $d$ are coprime
(Lemma~\ref{l:arith}(d)), the ideals $(y)$ and $(\sqd{P}{\mu})$ of $\OO$ are coprime. The different of $\OO$ over
$\Z$ is a multiple of $(y)$, by transitivity through $\Q(z)$, and it is the product of the different of $\OO$ over
$\OO_K$ and $(\sqd{P}{\mu})$, by transitivity through $\KP{P}$. Hence $(y)$ divides the different of $\OO$ over
$\OO_K$, and (I2), for the polynomial $\Cy{\Nn{P}{\mu}}$, gives both remaining assertions. (I2) applies because $z$
generates $\EP{P}{\mu}$ over $\KP{P}$: the root of unity $\rtu{P}{\mu}$ is a power of $z$, both being primitive of
order $\Nn{P}{\mu}$.
\end{proof}
\leanline{\lean{relspec}}

\begin{lemma}\label{l:dif}
Let $P=(N_0,d)$ be parameters and $\mu$ a natural number. Then $\dgen{P}{\mu}$ divides $\Nn{P}{\mu}$ in the ring of
integers of $\EP{P}{\mu}$, so $\dgen{P}{\mu}\ne0$; and $\Nn{P}{\mu}\,d$ lies in the different of that ring over $\Z$.
\end{lemma}
\begin{proof}
The first assertion is Theorem~C(a) for $z=\rti{P}{\mu}$. By Theorem~C(a) and transitivity the different over $\Z$
is generated by $\dgen{P}{\mu}\cdot\sqd{P}{\mu}$, which is the first half of Theorem~C(b); this product divides
$\Nn{P}{\mu}\,d$ by Lemma~\ref{l:ints}.
\end{proof}
\leanline{\lean{dgendvd}\sep\lean{dgenne}\sep\lean{difmul}\sep\lean{condmem}}

\begin{lemma}[A unit that makes the generator real]\label{l:real}
Let $P=(N_0,d)$ be parameters and $\mu$ a natural number. There is an integer $u_0$ of $\EP{P}{\mu}$ with
$\cj{u_0}\,u_0=1$ and $\cj{u_0\,\dgen{P}{\mu}}=u_0\,\dgen{P}{\mu}$.
\end{lemma}
\begin{proof}
The derivative of a product $\prod(X-x)$ of distinct linear factors at one of its roots is the product of the
differences of that root and the others. So $\dgen{P}{\mu}$ is the product of the differences
$\rti{P}{\mu}-x$ over the primitive roots of unity $x\ne\rti{P}{\mu}$ of order $\Nn{P}{\mu}$. Each of them is
$x=\rti{P}{\mu}^{\,l}$ with $l$ prime to $\Nn{P}{\mu}$, hence odd, $l=2i+1$, and
\[\rti{P}{\mu}-\rti{P}{\mu}^{\,2i+1}=\rti{P}{\mu}^{\,i+1}\bigl(\rti{P}{\mu}^{-i}-\rti{P}{\mu}^{\,i}\bigr).\]
The second factor changes sign under conjugation (Lemma~\ref{l:dict}(b)). As $4$ divides $\Nn{P}{\mu}$, the ring
of integers contains a square root $\mathrm{i}_\ast$ of $-1$, and $\cj{\mathrm{i}_\ast}=-\mathrm{i}_\ast$. Put
$u_x=\mathrm{i}_\ast\,\cj{\rti{P}{\mu}}^{\,i+1}$; then $\cj{u_x}\,u_x=1$ by Lemma~\ref{l:ints}, and
$u_x\,(\rti{P}{\mu}-x)=\mathrm{i}_\ast\bigl(\rti{P}{\mu}^{-i}-\rti{P}{\mu}^{\,i}\bigr)$ is fixed by conjugation.
The product $u_0$ of the $u_x$ has the two properties. (The formal proof is this induction over the factors.)
\end{proof}
\leanline{\lean{unitreal}}

\begin{definition}[The purely imaginary generator of the different]\label{d:varpi}
Let $P=(N_0,d)$ be parameters and $\mu$ a natural number. With an integer $u_0$ chosen as in Lemma~\ref{l:real},
$\vpi{P}{\mu}=u_0\cdot\dgen{P}{\mu}\cdot\sqd{P}{\mu}$.
\end{definition}
\leanline{\lean{vpi}}

\begin{lemma}\label{l:varpi}
Let $P=(N_0,d)$ be parameters, $\mu$ a natural number and $\tau\in\PlP{P}{\mu}$. Then
$\cj{\vpi{P}{\mu}}=-\vpi{P}{\mu}$, and $|\emb{\tau}(\vpi{P}{\mu})|=|\emb{\tau}(\dgen{P}{\mu})|\cdot\sqrt d$.
\end{lemma}
\begin{proof}
$u_0\,\dgen{P}{\mu}$ is fixed by conjugation and $\sqd{P}{\mu}$ changes sign. Further
$|\emb{\tau}(u_0)|^2=\emb{\tau}(\cj{u_0}\,u_0)=1$
and $|\emb{\tau}(\sqd{P}{\mu})|=\sqrt d$.
\end{proof}
\leanline{\lean{starvpi}\sep\lean{absvpi}}

\begin{proof}[Proof of Theorem C(b), (c)]
The first half of (b) was shown with Lemma~\ref{l:dif}; the second follows because $u_0$ is a unit. For (c), all
infinite places of $\EP{P}{\mu}$ are complex, so the square of the product in question is the absolute value of the
norm of $\vpi{P}{\mu}$ to $\Q$, that is the absolute norm of the different, that is the absolute value of the
discriminant, which is $\rdE{P}{\mu}$ to the power twice the number of infinite places. Take square roots.
\end{proof}
\leanline{\lean{difvpi}\sep\lean{prodvpi}}

\sect{Frobenius elements}{s:frob}
\begin{definition}[The action on finite places; the Frobenius element]\label{d:frob}
Let $P=(N_0,d)$ be parameters and $\mu$ a natural number, and let $\OO$ be the ring of integers of $\EP{P}{\mu}$.
For $\tau\in\PlP{P}{\mu}$ let $\galv{\tau}\colon\OO\to\OO$ be the restriction of $\galP{P}{\mu}(\tau)$ (its one explicit
argument in the formal text is $\tau$). The group $\PlP{P}{\mu}$ acts on the finite places of $\EP{P}{\mu}$ by
$\tau\cdot w=\{x\in\OO:\galv{\tau}(x)\in w\}$. For a natural number $p$,
\[\Frb{\mu}{p}=\bigl(\pin{P}{\mu}{\chiK{P}(p)\,p},\ \chiK{P}(p)\bigr)\in\PlP{P}{\mu}\times\Z,\]
and powers of such pairs are taken componentwise. (The parameters $P$ are an implicit argument of the formal
definition; the notation shows the two explicit ones, the exponent $\mu$ and $p$.)
\end{definition}
\leanline{\lean{galv}\sep\lean{Frb}\sep\lean{coegal}\sep\lean{memsmul}; the action is an instance without a name.}

\begin{setting}\label{set:fr}
$P=(N_0,d)$ are parameters, $\mu$ is a natural number, $\OO$ is the ring of integers of $\EP{P}{\mu}$, $p$ is a prime
number, $\tau\in\PlP{P}{\mu}$, and $(\mathfrak{c},e)$ is one of the two pairs $(\mathrm{id},1)$ and
$(\text{complex conjugation},-1)$, where $\mathfrak{c}$ is an automorphism of the ring $\OO$ and $e$ an integer. Further $R$ is
a commutative ring of characteristic $p$, $\psi\colon\OO\to R$ is a homomorphism of rings, and $n$ is a natural
number.
\end{setting}

\begin{lemma}[A Frobenius congruence is tested on two generators]\label{l:frobext}
In Setting~\ref{set:fr} assume $p\nmid\Nn{P}{\mu}$. Then $\psi(\galv{\tau}(x))=\psi(\mathfrak{c}(x))^{p^n}$ holds for all
$x\in\OO$ if and only if both
\[\psi\bigl(\galv{\tau}(\rti{P}{\mu})\bigr)=\psi\bigl(\galv{\tau'}(\rti{P}{\mu})\bigr)\quad\text{with }\tau'=\pin{P}{\mu}{e\,p^n}\]
and
\[\chiK{P}(p)^n\,e\,\psi(\sqd{P}{\mu})=\psi(\sqd{P}{\mu}).\]
\end{lemma}
\begin{proof}
Both $x\mapsto\psi(\galv{\tau}(x))$ and $x\mapsto\psi(\mathfrak{c}(x))^{p^n}$ are homomorphisms of rings, the second
because $R$ has characteristic $p$. By (I1) and Theorem~C(a) the ring $\OO$ is generated by $\rti{P}{\mu}$ and
$\czE{P}{\mu}$, so the two maps agree if and only if they agree at these two elements. At the root of unity:
$\mathfrak{c}$ raises it to the power $e$ (Lemma~\ref{l:dict}(b)), and $e\,p^n$ is prime to $N_0$ because
$p\nmid\Nn{P}{\mu}$, so $\mathfrak{c}(\rti{P}{\mu})^{p^n}=\galv{\tau'}(\rti{P}{\mu})$ by Theorem~A(b). At
$\czE{P}{\mu}$, which $\galv{\tau}$ fixes: here $d\equiv7\bmod8$ and $p$ is odd (Lemma~\ref{l:arith}(a) and the
hypothesis $p\nmid\Nn{P}{\mu}$), so (I1$'$) applies, with $\chi(p)=\chiK{P}(p)$, to the restriction $\psi'$ of
$\psi\circ\mathfrak{c}$ to the integers of $\KP{P}$, where $\psi'(\sqrt{-d})=e\,\psi(\sqd{P}{\mu})$, and to
$x=\psi(\czE{P}{\mu})$, where $2x-1=\psi(\sqd{P}{\mu})$. This turns the condition at $\czE{P}{\mu}$ into the second
equation.
\end{proof}
\leanline{\lean{frobext}\sep\lean{copNz}}

\begin{lemma}[Frobenius congruences and powers of the Frobenius element]\label{l:frobpow}
In Setting~\ref{set:fr} assume that $R$ is an integral domain and that $p\nmid\Nn{P}{\mu}\,d$. Then
$\psi(\galv{\tau}(x))=\psi(\mathfrak{c}(x))^{p^n}$ holds for all $x\in\OO$ if and only if $(\tau,e)=\Frb{\mu}{p}^{\,n}$ in
$\PlP{P}{\mu}\times\Z$.
\end{lemma}
\begin{proof}
Use Lemma~\ref{l:frobext}. The image of the root of unity in the domain $R$ is a root of the cyclotomic polynomial
of an order that is not zero in $R$, so its powers with exponents below that order are distinct; by Theorem~A(a)
the first condition says that $\tau$ and $\tau'$ have the same class, that is $\tau=\pin{P}{\mu}{e\,p^n}$. In the second
condition $\psi(\sqd{P}{\mu})$ can be cancelled, because its square $-d$ is not zero in $R$; it says
$\chiK{P}(p)^n\,e=1$ in $R$, hence $e=\chiK{P}(p)^n$, as both are $\pm1$ and $p\ne2$. Together with
Lemma~\ref{l:dict}(c) the two conditions say $\tau=\pin{P}{\mu}{\chiK{P}(p)\,p}^{\,n}$ and $e=\chiK{P}(p)^n$.
\end{proof}
\leanline{\lean{frobpow}}

\begin{proof}[Proof of Theorem B]
The residue ring $\OO/w$ is a finite field of characteristic $p$; let $\psi$ be the reduction. (a) If
$\tau\cdot w=\mathfrak{c}^{-1}(w)$, then $\galv{\tau}\circ\mathfrak{c}^{-1}$ maps $w$ into itself and induces an
endomorphism of the finite field, which is a power of the Frobenius map: for some $n$,
$\psi(\galv{\tau}(x))=\psi(\mathfrak{c}(x))^{p^n}$ for all $x$, and Lemma~\ref{l:frobpow} gives
$(\tau,e)=\Frb{\mu}{p}^{\,n}$. Conversely, such an $n$ gives that congruence by Lemma~\ref{l:frobpow}, and then
$x\in \tau\cdot w$ if and only if $\mathfrak{c}(x)^{p^n}\in w$, if and only if $\mathfrak{c}(x)\in w$.
(b) The norm of $w$ is $p^i$ for the residue degree $i$. By Lemma~\ref{l:frobpow}, for the unit place and the
first pair, $\Frb{\mu}{p}^{\,n}=1$ if and only if $x^{p^n}=x$ on the residue field, if and only if $i$ divides $n$.
So $i$ is the order of $\Frb{\mu}{p}$.
\end{proof}
\leanline{\lean{decomp}\sep\lean{qpow}}

\begin{lemma}[No ramification]\label{l:unram}
Let $P=(N_0,d)$ be parameters, $\mu$ a natural number, $p$ a prime number with $p\nmid\Nn{P}{\mu}\,d$ and $w$ a
finite place of $\EP{P}{\mu}$ with $p\in w$. Then $p\notin w^2$, and $p$ has valuation exactly $1$ at $w$.
\end{lemma}
\begin{proof}
If $p\in w^2$, then $w$ divides the different by (I3), which contains $\Nn{P}{\mu}\,d$ (Lemma~\ref{l:dif}); a natural
number in $w$ is divisible by $p$.
\end{proof}
\leanline{\lean{notsq}\sep\lean{valp}}

\begin{lemma}[Transitivity up to conjugation]\label{l:trans}
Let $P=(N_0,d)$ be parameters, $\mu$ a natural number, $p$ a prime number, and $w$, $w'$ finite places of
$\EP{P}{\mu}$ with $p\in w$ and $p\in w'$. Then there is $\tau\in\PlP{P}{\mu}$ with $\tau\cdot w=w'$ or
$\cj{\tau\cdot w}=w'$.
\end{lemma}
\begin{proof}
The Galois group of $\EP{P}{\mu}$ over $\Q$ permutes the primes above $p$ transitively: there is an automorphism
$\gamma$ with $\gamma(w')=w$, that is $w'=\{x:\gamma(x)\in w\}$. Let $\tau$ be the place of the embedding
$\emb{1}\circ\gamma$. The embedding $\emb{\tau}=\emb{1}\circ\galP{P}{\mu}(\tau)$ has the same place, so $\gamma$ is
$\galP{P}{\mu}(\tau)$, or $\galP{P}{\mu}(\tau)$ followed by complex conjugation, which commutes with $\galP{P}{\mu}(\tau)$
because $\EP{P}{\mu}$ is abelian over $\Q$ (Lemma~\ref{l:cyclo}). In the first case $w'=\tau\cdot w$ by the definition
of the action, which takes preimages; in the second
$w'=\{x:\cj{\galv{\tau}(x)}\in w\}=\{x:\galv{\tau}(\cj x)\in w\}=\cj{\tau\cdot w}$.
\end{proof}
\leanline{\lean{transit}}

\begin{lemma}[Primes inert in the quadratic field]\label{l:inert}
Let $P=(N_0,d)$ be parameters, $\mu$ a natural number and $p$ a prime number with $\chiK{P}(p)=-1$. Then
$p\nmid\Nn{P}{\mu}\,d$. Let moreover $w$ be a finite place of $\EP{P}{\mu}$ with $p\in w$, and $\tau\in\PlP{P}{\mu}$.
Then $\tau\cdot w=w$ if and only if $\tau=\pin{P}{\mu}{-p}^{\,n}$ for an even natural number $n$; $\tau\cdot w=\cj{w}$ if and
only if $\tau=\pin{P}{\mu}{-p}^{\,n}$ for an odd $n$; and the norm of $w$ is $p$ to the least common multiple of $2$
and the order of $\pin{P}{\mu}{-p}$ in $\PlP{P}{\mu}$.
\end{lemma}
\begin{proof}
The primes dividing $\Nn{P}{\mu}\,d$ have $\chiK{P}=1$ or $0$ (Lemma~\ref{l:arith}(b)). Here
$\Frb{\mu}{p}=(\pin{P}{\mu}{-p},-1)$, and Theorem~B applies with either pair.
\end{proof}
\leanline{\lean{inertnd}\sep\lean{inertfix}\sep\lean{inertstar}\sep\lean{inertq}}

\begin{lemma}[Numbers congruent to $1$ modulo the conductor]\label{l:split}
Let $P=(N_0,d)$ be parameters, $\mu$ a natural number and $p$ a natural number with
$p\equiv1\bmod\Nn{P}{\mu}\,d$. Then $\chiK{P}(p)=1$ and $\Frb{\mu}{p}=1$. If moreover $p$ is a prime number, then
$p\nmid\Nn{P}{\mu}\,d$; and for every finite place $w$ of $\EP{P}{\mu}$ with $p\in w$ and every $\tau\in\PlP{P}{\mu}$:
$\tau\cdot w\ne\cj{w}$; $\tau\cdot w=w$ if and only if $\tau=1$; and the norm of $w$ is $p$.
\end{lemma}
\begin{proof}
$p\equiv1$ modulo $d$ and modulo $\Nn{P}{\mu}$, so the symbol is $1$ and $\pin{P}{\mu}{p}$ is the unit element. A
prime dividing the modulus would divide $1$. The rest is Theorem~B with $\Frb{\mu}{p}=1$.
\end{proof}
\leanline{\lean{splitchi}\sep\lean{splitfr}\sep\lean{splitnd}\sep\lean{splitne}\sep\lean{splitfix}\sep\lean{splitq}}

\begin{lemma}[An integer of norm $p$ times a number prime to $p$]\label{l:norm}
Let $P=(N_0,d)$ be parameters, $\mu$ a natural number, $p$ a prime number with $p\equiv1\bmod\Nn{P}{\mu}\,d$, $w$ a
finite place of $\EP{P}{\mu}$ with $p\in w$, and $x_0$ an integer of $\EP{P}{\mu}$ such that $x_0\,\cj{x_0}=p\,n_0$ for a
natural number $n_0$ prime to $p$. Then $\cj{x_0}\in w$ if and only if $x_0\notin w$. If $x_0\in w$, then $x_0$ has valuation
exactly $1$ at $w$.
\end{lemma}
\begin{proof}
$x_0\,\cj x_0\in w$, so one of the two lies in $w$. If both did, then $p\,n_0\in w^2$ with $n_0\notin w$, so $p\in w^2$,
against Lemma~\ref{l:unram}. The valuations of $x_0$ and $\cj x_0$ at $w$ add up to that of $p\,n_0$, which is $1$.
\end{proof}
\leanline{\lean{starmem}\sep\lean{valnorm}}

\begin{proposition}[As many places above $p$ and $x_0$ as infinite places]\label{p:bij}
Let $P=(N_0,d)$ be parameters, $\mu$ a natural number, $p$ a prime number with $p\equiv1\bmod\Nn{P}{\mu}\,d$, and $x_0$
an integer of $\EP{P}{\mu}$ such that $x_0\,\cj{x_0}=p\,n_0$ for a natural number $n_0$ prime to $p$. Then there is a
bijection from $\PlP{P}{\mu}$ onto the set of finite places of $\EP{P}{\mu}$ that contain both $x_0$ and $p$.
\end{proposition}
\begin{proof}
Fix a finite place $w$ containing $p$. Send $\tau$ to $\tau\cdot w$ if $x_0\in \tau\cdot w$ and to $\cj{\tau\cdot w}$ otherwise.
Both places contain $p$, and exactly one of them contains $x_0$ (Lemma~\ref{l:norm}), so the map is well defined. If
$\tau'$ and $\tau$ have the same image, then $\tau'\cdot w$ and $\tau\cdot w=(\tau'^{-1}\tau)\cdot(\tau'\cdot w)$ are equal or conjugate;
Lemma~\ref{l:split} excludes the second and turns the first into $\tau'=\tau$. A place containing $x_0$ and $p$ is
$\tau\cdot w$ or $\cj{\tau\cdot w}$ for some $\tau$ (Lemma~\ref{l:trans}), and then it is the image of $\tau$.
\end{proof}
\leanline{\lean{bij}}

\sect{Weil data}{s:weil}
\begin{lemma}\label{l:valmap}
Let $P=(N_0,d)$ be parameters, $\mu$ a natural number, $\eta$ an injective endomorphism of the ring of integers of
$\EP{P}{\mu}$ and $w$ a finite place of $\EP{P}{\mu}$. Then $\eta(\dgen{P}{\mu})$ and $\dgen{P}{\mu}$ have the same
valuation at $w$.
\end{lemma}
\begin{proof}
$\eta(\dgen{P}{\mu})$ is the derivative of the cyclotomic polynomial at the primitive root $\eta(\rti{P}{\mu})$, and by
Theorem~C(a) it generates the same ideal as $\dgen{P}{\mu}$.
\end{proof}
\leanline{\lean{valmap}}

\begin{definition}[The idele of the different]\label{d:alpha}
Let $P=(N_0,d)$ be parameters and $\mu$ a natural number. $\alP{P}{\mu}$ is the finite idele of $\EP{P}{\mu}$ whose
component is $\dgen{P}{\mu}^{-1}$ at the finite places that contain $\czE{P}{\mu}$, and $1$ at all others.
\end{definition}
\leanline{\lean{alP}}

\begin{lemma}\label{l:alpha}
Let $P=(N_0,d)$ be parameters, $\mu$ a natural number and $w$ a finite place of $\EP{P}{\mu}$. Then the inverse of
$\alP{P}{\mu}$ and the product $\Nn{P}{\mu}\cdot\alP{P}{\mu}$ are integral at every finite place, and the component
at $w$ of $\alP{P}{\mu}\cdot\cj{\alP{P}{\mu}}\cdot\dgen{P}{\mu}$ is a unit.
\end{lemma}
\begin{proof}
The first two assertions hold because $\dgen{P}{\mu}$ is an integer that divides $\Nn{P}{\mu}$. For the third, the
component of the conjugate idele at $w$ is $\cj{\dgen{P}{\mu}}^{-1}$ if $\cj{\czE{P}{\mu}}=1-\czE{P}{\mu}$ lies in
$w$, and $1$ otherwise. Not both $\czE{P}{\mu}$ and $1-\czE{P}{\mu}$ lie in $w$. If exactly one does, the component
in question is $1$ or $\dgen{P}{\mu}/\cj{\dgen{P}{\mu}}$, a unit by Lemma~\ref{l:valmap} for complex conjugation.
If neither does, the component is $\dgen{P}{\mu}$, and $\dgen{P}{\mu}\notin w$: otherwise $\Nn{P}{\mu}\in w$, and as
$\Nn{P}{\mu}$ divides a power of $(d+1)/4=\czE{P}{\mu}(1-\czE{P}{\mu})$ by Lemma~\ref{l:arith}(b), one of the two
factors would lie in $w$.
\end{proof}
\leanline{\lean{alinv}\sep\lean{valpha}\sep\lean{Nal}}

\statement{ana}{\textbf{Inputs.}\enspace In (I4) to (I6), $F$ is a CM field with a chosen embedding $\emb{\tau}$ at
every infinite place $\tau$. (I4), unit~106: \emph{Weil data} for $F$ consist of a primitive root of unity of order
$5040$ in $F$, a \emph{splitting character} (a continuous unitary character $\upsilon$ of the ideles of $F$,
trivial on $F^\times$, with $\upsilon(\cj x)=\upsilon(x)^{-1}$ for every idele $x$, and with
$\upsilon(x)=\prod_\tau x_\tau/|x_\tau|$ for every idele $x$ whose finite components are all local units), an element
$\mathbf{j}\in F$, and a finite idele, such that: $\cj{\mathbf{j}}=-\mathbf{j}$; the imaginary part of $\emb{\tau}(\mathbf{j})$ is
positive for every infinite place $\tau$; the product of $\mathbf{j}$, the idele and its conjugate generates the inverse
of the different over $\Z$ at every finite place; $(1-\mathbf{j}^{-1})/2$ is an algebraic integer; and $\upsilon$ takes
the value $-1$ at a uniformiser of every finite place $w$ with $\cj w=w$.
(I5), unit~104: let $F$ contain a primitive root of unity of order $5040$ (the formal text calls such a CM field
\emph{stably unramified}: then $F$ contains a square root of $-1$ and $(1+\sqrt{-7})/2$, so it is unramified over
its maximal totally real subfield at every finite place). If every root of unity $y$ of $F$ is $y=u/\cj u$ for a
unit $u$ of the ring of integers with $\prod_\tau\emb{\tau}(u)=1$, then $F$ has a splitting character. (I6), unit~97: every splitting character takes the value $-1$
at a uniformiser of every finite place $w$ with $\cj w=w$.}

\begin{proposition}[Weil data for the CM field]\label{p:weil}
Let $P=(N_0,d)$ be parameters and $\mu$ a natural number. The field $\EP{P}{\mu}$, with the embeddings of
Definition~\ref{d:group}, has Weil data in the sense of (I4), with $\mathbf{j}=-1/\sqd{P}{\mu}$ and the idele
$\alP{P}{\mu}$.
\end{proposition}
\begin{proof}
As $5040=2^4\cdot3^2\cdot5\cdot7$ divides $\Nn{P}{\mu}$ ($2^4\cdot5\cdot7$ divides $18480$, hence $N_0$, and
$3^2$ divides $3^5$, hence $\twd{\mu}$), a power of $\rti{P}{\mu}$ has order $5040$: this is the first datum of
(I4) and the first hypothesis of (I5). For the second hypothesis of (I5), let $y$ be a root of unity;
$y=\rti{P}{\mu}^{\,n}$ by Lemma~\ref{l:roots}(a). Put $i=\Nn{P}{\mu}/4$, a natural number, and
\[u=\rti{P}{\mu}^{\,i}\,\bigl(\rti{P}{\mu}-1\bigr).\]
Its norm to $\KP{P}$ is $1$: the norm of $\rti{P}{\mu}$ is $1$, and the norm of $\rti{P}{\mu}-1$ is the value
$\Cy{\Nn{P}{\mu}}(1)$, by Lemma~\ref{l:irr} and because the degree is even; that value is $1$ since $\Nn{P}{\mu}$ is
not a power of a prime. So $u$ is a unit, and $\prod_\tau\emb{\tau}(u)=1$ by Lemma~\ref{l:roots}(b). Further
$\cj u=\rti{P}{\mu}^{-i}(\rti{P}{\mu}^{-1}-1)=-\rti{P}{\mu}^{-i-1}(\rti{P}{\mu}-1)$, and $\rti{P}{\mu}^{\,2i}=-1$, so
$\rti{P}{\mu}\,\cj u=u$; hence $y\,\cj{u^n}=u^n$, and (I5) gives a splitting character. Now $\mathbf{j}$: its conjugate
is $-\mathbf{j}$ by Lemma~\ref{l:ints}; $\emb{\tau}(\mathbf{j})=\mathrm{i}/\sqrt d$ for every $\tau$, because $\emb{\tau}$ is the
identity on $\KP{P}$; and $(1-\mathbf{j}^{-1})/2=\czE{P}{\mu}$. By Lemma~\ref{l:alpha} and Theorem~C(b),
$\mathbf{j}\,\alP{P}{\mu}\,\cj{\alP{P}{\mu}}$ has at every finite place the valuation of the inverse of
$\dgen{P}{\mu}\cdot\sqd{P}{\mu}$, a generator of the different. The last condition is (I6).
\end{proof}
\leanline{\lean{weil}}

\sect{Notes on the formal text}{s:notes}
\subhead{Names and arguments.} The formal text writes $N_0$ and $d$ as the two components of $P$, and every other
object of this note as a function of $P$ and the exponent; the note keeps all explicit arguments. Three notations
show less than the full list of arguments, and say so where they are defined: the Frobenius element and the
permutation of places have $P$ as an implicit argument, and the restriction $\galv{\tau}$ has both $P$ and $\mu$
implicit. The class map takes $P$ and a proof of a divisibility; the note writes the divisibility. The formal text
calls the exponent \texttt{s} and the root of unity \texttt{ζ}; this write-up keeps those letters for the complex
variable and for Riemann's function.

\subhead{Signatures that the reader of the unit derived.} For eleven declarations of this note the list of arguments
in the dossier is derived from Lean's rule for section variables and not printed by Lean: Theorem~B(a), the
statements of Lemmas~\ref{l:trans}, \ref{l:norm} and Proposition~\ref{p:bij}, the three statements on places in
Lemma~\ref{l:split}, the two equivalences of Lemma~\ref{l:inert}, and Theorem~C(a). I read each of them in the Lean
text and agree: where the statement binds a variable itself (the place $\tau$ in Lemma~\ref{l:trans}, the exponent in
Theorem~B(a) and Lemma~\ref{l:inert}, the element $y$ of Theorem~C(a), the cofactor $n_0$ of Lemma~\ref{l:norm}), the
section variable of the same name is not an argument; and the hypothesis $p\in w$ is an argument of the three
statements of Lemma~\ref{l:split} that mention $w$, because the text includes it again after having omitted it.
This follows the documented rule; only Lean's own output can confirm it, and nothing was compiled here.

\subhead{The embedding into the cyclotomic field.} The census of the workspace records one explicit argument for
the embedding of Lemma~\ref{l:cyclo}, an element of $\EP{P}{\mu}$. In the Lean text the declaration is a
homomorphism of $\Q$-algebras with the two implicit arguments $P$ and $\mu$ and no explicit one; the arrow of the
homomorphism was read as an argument. The note accounts for the entry by naming the field.

\subhead{Where the formal statements differ from the natural ones.} Theorem~A(a) is stated in Lean with the class map for the divisor $\Nn{P}{\mu}$ of itself, which is $\cycP{P}{\mu}$. Lemma~\ref{l:arith}(b) is stated in Lean with
the quotient $(d+1)/4$ of natural numbers; it is exact. In Theorem~B and Lemmas~\ref{l:frobext}, \ref{l:frobpow} the
pair $(\mathfrak{c},e)$ is a ring automorphism and an integer with the hypothesis that they form one of the two pairs; the
group $\PlP{P}{\mu}\times\Z$ has the integers as a multiplicative monoid, so the order in Theorem~B(b) is the order
in a monoid. Valuations are multiplicative in Lean: ``valuation exactly $1$'' is there the value, in the value group of
the discrete valuation, that Lean writes \texttt{.exp (-1)}, the value of a uniformiser (not the real number
$\mathrm{e}^{-1}$). The census of the workspace cuts the statement of Theorem~C(a) at the word \texttt{let}; the statement
runs two lines further, as the reader's corrected list of statement spans records and as the Lean file shows, and
for this reason it is not among the statements quoted in Appendix~A. In
Lemma~\ref{l:split} the first two assertions do not assume that $p$ is prime; the formal text removes that
hypothesis explicitly. Lemma~\ref{l:dict}(c) and (e) hold for all integers $\nu$, by the convention that an integer
not prime to the order names the unit place. Proposition~\ref{p:weil} is an instance: it provides the data rather
than asserting their existence, and it depends on the choices of Lemma~\ref{l:real} and of (I5). The formal
definition behind (I4) lists the root of unity of order $5040$ as the property that (I5) assumes, and it writes
$\mathbf{j}$ as \texttt{j}; in the handout the plain letter \emph{j} is the lift index, and a homomorphism of rings
from $\CC$ to a non-archimedean field has its own letter, so the element of the Weil data is written in bold. The first hypothesis of (I5) cannot be dropped. The field $\Q(\sqrt3,\sqrt{-2})$ has the roots of unity
$\pm1$, and $u=(\sqrt{-2}+\sqrt{-6})/2$ is a unit in it with $u/\cj u=-1$ and, for the two embeddings that send
$\sqrt{-2}$ to the upper half-plane, with $\prod_\tau\emb{\tau}(u)=1$. But a splitting character of a CM field restricts
to the ideles of the maximal totally real subfield as the quadratic character of the extension, and its values
on the ideles with unit finite components show that this character is unramified at every finite place; the
field $\Q(\sqrt3,\sqrt{-2})$ is ramified over $\Q(\sqrt3)$ above~$2$. (This example is due to the referee of the
note and was checked by hand; it is not in the formal text.)

\subhead{What is used where.} Inside the unit, the note \emph{Admissible data} uses Theorem~A, Lemmas~\ref{l:arith}
and~\ref{l:dict}, Lemmas~\ref{l:unram} and~\ref{l:norm}, and Proposition~\ref{p:bij}. The remaining results serve later units: Lemma~\ref{l:perm} the
conjugates of algebraic parts, Lemmas~\ref{l:inert} and~\ref{l:powd} the local analysis at the primes that do not
split in $\KP{P}$, Theorem~C and Proposition~\ref{p:weil} the Weil representation. This was seen from the places of use
in the source, not from a kernel export.

\subhead{How this note was written.} From the dossier of the unit: the texts of its members and the structured
proofs of the declarations of the two files treated here. I read both Lean files completely and compared every
statement of this note with the Lean statement, argument by argument, with the lists of arguments of the census
and of the dossier. The proofs follow the structured proofs; those of Lemma~\ref{l:arith}(d),
Lemma~\ref{l:dict}(c) and Theorem~A(c), (d), which the formal text leaves to automation, are mine. The Inputs were
taken from the dossier's reading of the library and from the library's own comments; I opened the definition
(I4) and the statements of (I0), (I1$'$) and (I5), and none of the proofs behind (I0) to (I6). In the revision of
26~November 2026, after a referee's report, the hypothesis of (I5) on the root of unity of order $5040$ was
restored from the formal statement, and (I0) and (I1$'$) were moved from the proofs into the box of Inputs. A second reading of the unit by the reading agent arrived while this
note was written; its three corrections that touch this note (the polynomial that defines the field, the direction
of the automorphism in Lemma~\ref{l:trans}, the value of a uniformiser) were checked against the Lean text and
adopted. The program check of the unit's dossier covered \UnitSteps{} leaf steps
of \UnitProofs{} proofs.

\head{Appendix A. Concordance with the Lean artifact}
The formal names of the objects: \texttt{P.N₀}, \texttt{P.d} are $N_0$, $d$; \texttt{towerDeg} applied to $\mu$ is $\twd{\mu}$; the exponent \texttt{s} is $\mu$; \texttt{P.ζ s} and \texttt{P.zetaInt s} are the root of unity;
\texttt{P.Φ s} is the group of places and \texttt{Φ v} the embedding $\emb{\tau}$; \texttt{star} is complex conjugation.
Three statements as the artifact has them:
\quoted{par}
\quoted{prodvpi}
\quoted{decomp}
{\fontsize{8.6}{11.6}\selectfont\renewcommand{\leanfont}{\ttfamily\fontsize{7.7}{10}\selectfont\color{leancol}}
\setlength{\LTleft}{0pt}\setlength{\LTright}{\fill}\setlength{\tabcolsep}{0pt}
\begin{longtable}{@{}p{27mm}p{62mm}p{40mm}r@{}}
\textbf{In this note} & \textbf{Lean declaration} & \textbf{File} & \textbf{Lines}\\\midrule\endhead
\ConcordanceRows
\end{longtable}}

\head{Notation}
\begingroup\fontsize{8.6}{10.9}\selectfont\raggedright
\setlength{\columnsep}{9mm}\setlength{\parskip}{0pt}
\begin{multicols}{2}
\nt{$P=(N_0,d)$}{parameters (Definition~\ref{d:par})}
\nt{$\mu$}{the exponent, a natural number}
\nt{$\twd{\mu}$, $\Nn{P}{\mu}$}{the tower degree; the order of the root of unity}
\nt{$\Mom{P}{\mu}$}{the number $(N_0d)^{\mu+6}$}
\nt{$\KP{P}$, $\chiK{P}$}{$\Q(\sqrt{-d})$ and its quadratic character}
\nt{$\EP{P}{\mu}$, $\rtu{P}{\mu}$}{the CM field and its chosen root of unity}
\nt{$\rti{P}{\mu}$}{the same root of unity as an element of $\OO$ (Definition~\ref{d:ints})}
\nt{$\PlP{P}{\mu}$}{the group of infinite places of $\EP{P}{\mu}$}
\nt{$\OmP{P}{\mu}$, $\rdE{P}{\mu}$}{the period of Definition~\ref{d:E}; the root discriminant; $\mathrm{cs}(d)$ as there}
\nt{$\galP{P}{\mu}$, $\cycP{P}{\mu}$}{from places to automorphisms and to classes}
\nt{$\pin{P}{\mu}{\nu}$}{the place named by the integer $\nu$}
\nt{$\pcl{P}{m'\mid\Nn{P}{\mu}}$}{the class of a place modulo a divisor $m'$ of $\Nn{P}{\mu}$}
\nt{$\emb{\tau}$, $\emb{1}$}{the chosen embedding at the place $\tau$, at the unit place}
\nt{$\prm{\mu}{\sigma}$}{the permutation of places by $\sigma\in\Aut(\CC)$}
\nt{$\Cy{M}$}{the cyclotomic polynomial of order $M$ (the letter $\Psi$ is local to this note); $X$ its variable; $Q$ a cofactor}
\nt{$\OO$, $\OO_K$, $\OO_F$}{rings of integers}
\nt{$\czE{P}{\mu}$, $\sqd{P}{\mu}$}{$(1+\sqrt{-d})/2$ and $\sqrt{-d}$ in $\OO$}
\nt{$\dgen{P}{\mu}$, $\vpi{P}{\mu}$}{generators of the different over $\KP{P}$ and over $\Q$}
\nt{$\galv{\tau}$}{the restriction of $\galP{P}{\mu}(\tau)$ to $\OO$}
\nt{$\tau\cdot w$, $\cj{w}$}{the action of a place $\tau$ on a finite place $w$; the conjugate place}
\nt{$\Frb{\mu}{p}$}{the Frobenius element at $p$}
\nt{$\alP{P}{\mu}$}{the idele of Definition~\ref{d:alpha}}
\nt{$p$}{a prime number, or a natural number where said}
\nt{$\tau,\tau',\tau^\ast$}{infinite places}
\nt{$w,w'$}{finite places}
\nt{$\mathfrak{c},e$}{an automorphism of $\OO$ and a sign (Setting~\ref{set:fr})}
\nt{$R$, $\psi$, $\psi'$}{a ring; homomorphisms and embeddings (local to this note)}
\nt{$\chi(p)$}{in (I1$'$): the Legendre symbol of $p$ modulo $d$}
\nt{$\varphi$, $\pi$, $L$}{Euler's function; the circle constant; in ``$L$-function''}
\nt{$F,F'$}{auxiliary number fields}
\nt{$x$}{In Sections~\ref{s:field} to~\ref{s:weil} also an element of a field or of a ring}
\nt{$y,z,x_0$}{elements of fields and rings}
\nt{$n,i,l,n_0,o$}{natural numbers and integers: exponents, orders, degrees, a cofactor}
\nt{$\nu$}{In Section~\ref{s:setting} and after it also an integer naming a place, never a character of the ideles (Convention~0.6)}
\nt{$\sigma$, $\gamma$}{an automorphism of $\CC$; an element of a Galois group}
\nt{$\mathbf{j}$}{the purely imaginary element of the Weil data}
\nt{$\upsilon$}{a splitting character}
\nt{$\eta$}{an endomorphism of $\OO$}
\nt{$u_0$, $u_x$}{the unit of Lemma~\ref{l:real} and its factors}
\nt{$u$}{a unit, in (I5) and in the proof of Proposition~\ref{p:weil}}
\end{multicols}
\endgroup

\head{References}
\begingroup\fontsize{9.4}{12.6}\selectfont
\begin{thebibliography}{9}
\bibitem{conrad} K. Conrad, \emph{The different ideal}, expository notes on the author's web page. Theorem~4.13
and its proof, as recorded by the reading agent for (I3); not opened for this note.
\bibitem{neukirch} J. Neukirch, \emph{Algebraic Number Theory}, 1999. As recorded in the catalogue of the
write-up, places not confirmed.
\bibitem{serre} J.-P. Serre, \emph{Local Fields}. As recorded in the catalogue of the write-up, place not
confirmed.
\end{thebibliography}
\endgroup

\end{document}
